%% 3-6 2本の直線が交わるところ：連立方程式（第12回：組 V3-3。ページ制限なし。偶数ページ）  % noruby
%% 2-16 の約束を受ける：「2本の線が同じかたむきで平行になり、いつまでも交わりません（3-6）」  % noruby
%%   「電気の回路では、わからない電流をいくつも文字にして、方程式を何本もいっしょに解きます（連立方程式、3-6）」  % noruby
\makeatletter
% 式を、まん中に1行で置く  \bkthrgline{…}  % noruby
\newcommand\bkthrgline[1]{\par\addvspace{3pt}{\centering #1\par}\addvspace{3pt}}
% 説明の段の見出し（丸数字＋太字の1行）  \bkthrgstep{番号}{見出し}  % noruby
\newcommand\bkthrgstep[2]{\par\addvspace{6pt}\Needspace{4\baselineskip}\noindent
  {\bkcircled[bkOrange]{#1}\enspace\sffamily\bfseries\color{bkOrange!75!black}#2}\par\nobreak\vspace{2pt}}
% みかん  \bkthrgmikan{(x,y)}{半径}  % noruby
\newcommand\bkthrgmikan[2]{\begin{scope}[shift={#1}]
  \fill[bkOrange!85] (0,0) circle (#2);
  \fill[bkLeaf] (0,{#2*0.85}) ellipse ({#2*0.35} and {#2*0.2});
  \fill[white, opacity=0.35] ({-#2*0.4},{#2*0.25}) circle ({#2*0.2});
\end{scope}}
% かい物のトレー  \bkthrgtray{(左下)}{りんごの数}{みかんの数}{ねだん}  % noruby
\newcommand\bkthrgtray[4]{\begin{scope}[shift={#1}]
  \fill[bkGray!20, rounded corners=1.5mm] (0,0) rectangle (44,12);
  \draw[bkGray!80, line width=0.6pt, rounded corners=1.5mm] (0,0) rectangle (44,12);
  \foreach \k in {1,...,#2} {\bkapple{({1.5+\k*6.5},6)}{2.7}{bkAppleRed}}
  \foreach \k in {1,...,#3} {\bkthrgmikan{({#2*6.5+3.5+\k*6.5},5.6)}{2.6}}
  \node[font=\sffamily\bfseries\small, anchor=west] at (45.5,6) {#4\R{円}{えん}};
\end{scope}}
% 回路（電源2つ・抵抗3つ）  \bkthrgcircuit{左の電源}{左の抵抗}{まん中の抵抗}{右の抵抗}{右の電源}  % noruby
%   左のえだ x=0、まん中 x=40、右 x=80。上の線 y=44、下の線 y=0。電源は＋（長い線）が上  % noruby
\newcommand\bkthrgbatt[1]{\begin{scope}[shift={#1}]
  \draw[bkInk, line width=0.9pt] (0,-12) -- (0,-1.5) (0,1.5) -- (0,11);
  \draw[bkInk, line width=0.9pt] (-4.2,1.5) -- (4.2,1.5);
  \draw[bkInk, line width=2.4pt] (-2,-1.5) -- (2,-1.5);
\end{scope}}
\newcommand\bkthrgres[1]{\begin{scope}[shift={#1}]
  \filldraw[fill=white, draw=bkInk, line width=0.9pt] (-2,-5) rectangle (2,5);
\end{scope}}
\newcommand\bkthrgcircuit[5]{%
  \draw[bkInk, line width=0.9pt] (0,23) -- (0,44) -- (80,44) -- (80,23);
  \draw[bkInk, line width=0.9pt] (0,0) -- (80,0);
  \draw[bkInk, line width=0.9pt] (40,44) -- (40,27) (40,17) -- (40,0);
  \bkthrgbatt{(0,12)}\bkthrgbatt{(80,12)}
  \bkthrgres{(0,28)}\bkthrgres{(80,28)}\bkthrgres{(40,22)}
  \node[font=\sffamily\scriptsize, anchor=south west, inner sep=0.5pt] at (4.6,13.6) {＋};
  \node[font=\sffamily\scriptsize, anchor=south east, inner sep=0.5pt] at (75.4,13.6) {＋};
  \node[font=\sffamily\small, anchor=east] at (-5.5,12) {#1};
  \node[font=\sffamily\small, anchor=east] at (-3.8,28) {#2};
  \node[font=\sffamily\small, anchor=west] at (43.8,22) {#3};
  \node[font=\sffamily\small, anchor=west] at (83.8,28) {#4};
  \node[font=\sffamily\small, anchor=west] at (85.5,12) {#5};
  \draw[bkRed, line width=1.1pt, -{Stealth[length=4pt]}] (6,44) -- (16,44);
  \node[font=\small, text=bkRed!80!black, anchor=south] at (11,45.3) {$I_1$};
  \draw[bkRed, line width=1.1pt, -{Stealth[length=4pt]}] (74,44) -- (64,44);
  \node[font=\small, text=bkRed!80!black, anchor=south] at (69,45.3) {$I_2$};
  \draw[bkRed, line width=1.1pt, -{Stealth[length=4pt]}] (40,41.5) -- (40,33);
  \node[font=\small, text=bkRed!80!black, anchor=west] at (41.6,37.5) {$I_1+I_2$};}
% 係数の表（2行3列）  \bkthrgtab{(左上)}{a}{b}{c}{d}{e}{f}  % noruby
\newcommand\bkthrgtab[7]{\begin{scope}[shift={#1}]
  \fill[white] (0,-12) rectangle (27,0);
  \draw[bkLine!80!black, line width=0.5pt] (0,-6) -- (27,-6) (9,0) -- (9,-12);
  \draw[bkNavy!70, line width=1pt] (18,0.6) -- (18,-12.6);
  \draw[bkInk!70, line width=0.7pt] (0,0) rectangle (27,-12);
  \node[font=\small] at (4.5,-3) {$#2$};\node[font=\small] at (13.5,-3) {$#3$};\node[font=\small] at (22.5,-3) {$#4$};
  \node[font=\small] at (4.5,-9) {$#5$};\node[font=\small] at (13.5,-9) {$#6$};\node[font=\small] at (22.5,-9) {$#7$};
\end{scope}}
\makeatother

\bkchapter{3}{6}{2\R{本}{ほん}の\R{直線}{ちょく|せん}が\R{交}{まじ}わるところ：\\\R{連立方程式}{れん|りつ|ほう|てい|しき}}{\bklevel{2}}{2\R{本}{ほん}の\R{直線}{ちょく|せん}が\R{交}{まじ}わるところ：\R{連立方程式}{れん|りつ|ほう|てい|しき}}
\bkmichishirube{2-13,2-14,2-16,2-18,3-5}{3-10,3-16,7-1,9-10}

\begin{goal}
\begin{itemize}[label={\color{\bkTT}●}, leftmargin=1.6\zw, itemsep=0pt, topsep=1pt]
\item \bklevel{2}\enspace 2\R{本}{ほん}の\R{直線}{ちょく|せん}の\R{交}{まじ}わる\R{点}{てん}が、2つの\R{式}{しき}を\R{同時}{どう|じ}にみたす $(x,\ y)$ であることがわかる。
\item \bklevel{2}\enspace \R{代入法}{だい|にゅう|ほう}と\R{加減法}{か|げん|ほう}で\R{連立方程式}{れん|りつ|ほう|てい|しき}を\R{解}{と}き、\R{代入}{だい|にゅう}して\R{確}{たし}かめられる。
\item \bklevel{2}\enspace \R{解}{かい}がない・\R{無数}{む|すう}にあるわけを\R{説明}{せつ|めい}し、\R{出会}{で|あ}い・\R{買}{か}い\R{物}{もの}・\R{回路}{かい|ろ}の\R{問題}{もん|だい}を\R{解}{と}ける。
\end{itemize}
\end{goal}

\begin{hatena}
\R{家}{いえ}と\R{駅}{えき}は\qty{1500}{m}はなれています。ミナさんは\R{家}{いえ}から\R{駅}{えき}へ、\R{分速}{ふん|そく}\qty{60}{m}で\R{歩}{ある}きます。ソラさんは\R{同}{おな}じときに\R{駅}{えき}を\R{出}{で}て、\R{家}{いえ}へ\R{向}{む}かって\R{分速}{ふん|そく}\qty{90}{m}で\R{歩}{ある}きます。2\R{人}{にん}が\R{出会}{で|あ}うのは\R{何分後}{なん|ぷん|ご}、\R{家}{いえ}から\R{何}{なん}mの\R{所}{ところ}かな？　2\R{人}{にん}の\R{位置}{い|ち}と\R{時間}{じ|かん}のグラフをかいたら、\R{出会}{で|あ}う\R{所}{ところ}はグラフのどこに\R{現}{あらわ}れるでしょう。
\begin{bkfigure}
\begin{tikzpicture}[x=1mm, y=1mm]
  % 1500 m を 93.75 mm でかく（1 m を 1/16 mm）
  \draw[bkGray!60, line width=3pt] (0,0) -- (93.75,0);
  \foreach \m in {0,300,...,1500} {\pgfmathsetmacro\xx{\m/16}
    \draw[bkInk, line width=0.7pt] (\xx,-1.6) -- (\xx,1.6);
    \node[bk tick, below=2mm] at (\xx,0) {\m};}
  \node[bk tick, anchor=north west] at (99,-2.5) {m};
  % 家
  \fill[bkOrange!70] (-8,2.5) rectangle (-2,8); \fill[bkRed!70] (-9.5,8) -- (-5,12) -- (-0.5,8) -- cycle;
  \node[bk label, anchor=south] at (-5,13) {\R{家}{いえ}};
  % 駅
  \fill[bkBlue!50] (97,2.5) rectangle (109,11);
  \fill[bkNavy!80] (96,11) rectangle (110,13.5);
  \node[bk label, anchor=south] at (103,14.5) {\R{駅}{えき}};
  \bkkid{(5,2.5)}{12}{bkRoute2}{bkHair}
  \draw[bk arrow, bkRoute2] (11,6.5) -- (21,6.5);
  \node[font=\sffamily\footnotesize, text=bkRoute2!80!black, anchor=south west, inner sep=1pt] at (11,8.3) {ミナ\enspace\qty{60}{m/min}};
  \bkkid{(89,2.5)}{12}{bkRoute4}{bkHair}
  \draw[bk arrow, bkRoute4] (83,6.5) -- (73,6.5);
  \node[font=\sffamily\footnotesize, text=bkRoute4!80!black, anchor=south east, inner sep=1pt] at (83,8.3) {ソラ\enspace\qty{90}{m/min}};
  \node[font=\sffamily\bfseries\Large, text=bkYellow!80!black] at (47,7) {?};
\end{tikzpicture}
\bkfigcaption{2\R{人}{にん}は\R{同}{おな}じときに\R{出発}{しゅっ|ぱつ}。どこで\R{出会}{で|あ}う？}
\end{bkfigure}
\end{hatena}

\noindent 2-16では、\R{兄}{あに}が\R{弟}{おとうと}に\R{追}{お}いつく\R{時間}{じ|かん}を、1\R{次方程式}{じ|ほう|てい|しき}で\R{求}{もと}めました。そのとき、2\R{人}{にん}のグラフの\R{交}{まじ}わる\R{点}{てん}が\R{答}{こた}えになっていました。この\R{章}{しょう}では、「2\R{本}{ほん}の\R{直線}{ちょく|せん}の\R{交}{まじ}わる\R{点}{てん}」をしっかり\R{考}{かんが}えます。わからない\R{数}{かず}が2つ（\R{何分後}{なん|ぷん|ご}と、どこで）あるとき、2つの\R{式}{しき}をいっしょに\R{使}{つか}って\R{解}{と}く\R{方法}{ほう|ほう}を\R{学}{まな}びます。

%% ---------------- やってみよう・思い出そう ----------------
\begin{yattemiyou}[\R{表}{ひょう}とグラフで、\R{出会}{で|あ}う\R{所}{ところ}をさがそう]
\begin{enumerate}[label={\bkcircled[bkGreen]{\arabic*}}, leftmargin=2.4\zw, itemsep=1pt, topsep=0pt]
\item ミナさんは\R{家}{いえ}から1\R{分}{ぷん}に\qty{60}{m}\R{進}{すす}みます。ソラさんは\qty{1500}{m}の\R{所}{ところ}から、1\R{分}{ぷん}に\qty{90}{m}ずつ\R{家}{いえ}に\R{近}{ちか}づきます。\R{下}{した}の\R{表}{ひょう}をノートに\R{写}{うつ}し、2\R{分}{ふん}ごとに、\R{家}{いえ}からの\R{位置}{い|ち}をうめよう。
\item \R{横}{よこ}に\R{時間}{じ|かん}、たてに\R{家}{いえ}からの\R{位置}{い|ち}をとって、2\R{人}{にん}の\R{点}{てん}を\R{色}{いろ}を\R{変}{か}えてかき、それぞれ\R{直線}{ちょく|せん}で\R{結}{むす}ぼう。
\item 2\R{本}{ほん}の\R{直線}{ちょく|せん}が\R{交}{まじ}わるのは、どこかな？　そのとき2\R{人}{にん}の\R{位置}{い|ち}はどうなっているかな？
\end{enumerate}
\begin{bkcenter}
\begin{bktable}{colspec={Q[l,m,wd=34mm] *{7}{Q[c,m,wd=9.4mm]}}, rowsep=1pt, colsep=3pt}
\R{時間}{じ|かん}（\R{分}{ふん}） & 0 & 2 & 4 & 6 & 8 & 10 & 12 \\
ミナの\R{位置}{い|ち}（m） & 0 & 120 & 240 & & & & \\
ソラの\R{位置}{い|ち}（m） & 1500 & 1320 & 1140 & & & & \\
\end{bktable}
\end{bkcenter}
\end{yattemiyou}

\begin{omoidasou}
\bkrecall{2-13}{\R{文字}{も|じ}に\R{式}{しき}を\R{入}{い}れるときは、かっこごと\R{入}{い}れる。$y=x+1$ なら $3y=3(x+1)$。}
\bkrecall{2-14}{\R{等式}{とう|しき}は\R{天}{てん}びん。\R{両辺}{りょう|へん}に\R{同}{おな}じ\R{数}{かず}をたしても、ひいても、かけても、\R{等式}{とう|しき}は\R{成}{な}り\R{立}{た}つ。}
\bkrecall{2-16}{1\R{次方程式}{じ|ほう|てい|しき}は、\R{移項}{い|こう}して\R{同類項}{どう|るい|こう}をまとめ、$x$ の\R{係数}{けい|すう}でわって\R{解}{と}く。$2t=10+2t$ には\R{解}{かい}がない。}
\bkrecall{2-18}{\R{表計算}{ひょう|けい|さん}では、\R{式}{しき}を\R{下}{した}へコピーすると、\R{行}{ぎょう}の\R{番号}{ばん|ごう}もずれる。}
\bkrecall{3-5}{1\R{次関数}{じ|かん|すう} $y=ax+b$ のグラフは\R{直線}{ちょく|せん}。$a$ が\R{傾}{かたむ}き、$b$ が\R{切片}{せっ|ぺん}。}
\end{omoidasou}

\noindent やってみようの\R{表}{ひょう}は、ミナさんが 360、480、600、720、ソラさんが 960、780、600、420 と\R{続}{つづ}きます。10\R{分後}{ぷん|ご}に、2\R{人}{にん}とも\qty{600}{m}の\R{所}{ところ}にいます。\R{同}{おな}じ\R{時間}{じ|かん}に\R{同}{おな}じ\R{場所}{ば|しょ}にいるので、ここで\R{出会}{で|あ}います。

\bkfukidashi{mina}{\R{表}{ひょう}で\R{見}{み}つかったけど、もし\R{出会}{で|あ}うのが9.4\R{分後}{ふん|ご}みたいな\R{半}{はん}ぱな\R{時間}{じ|かん}だったら、\R{表}{ひょう}では\R{見}{み}つけにくいね。}

%% ---------------- 数学の道具 ----------------
\bksection{doguu}{\R{数学}{すう|がく}の\R{道具}{どう|ぐ}}

\bksubsection{1}{\R{交}{まじ}わる\R{点}{てん}は、2つの\R{式}{しき}を\R{同時}{どう|じ}にみたす}

\leavevmode\label{bk:3-6:kouten}\R{出発}{しゅっ|ぱつ}してからの\R{時間}{じ|かん}を $x$ \R{分}{ふん}、\R{家}{いえ}からの\R{位置}{い|ち}を $y$〔m〕とします。2\R{人}{にん}の\R{位置}{い|ち}を\R{式}{しき}にしましょう。

\textbf{ミナさん}\quad \R{家}{いえ}（0の\R{所}{ところ}）から1\R{分}{ぷん}に\qty{60}{m}ずつ\R{進}{すす}むので、$y=60x$ です（\R{道}{みち}のり＝\R{速}{はや}さ×\R{時間}{じ|かん}。2-3）。

\textbf{ソラさん}\quad はじめは\qty{1500}{m}の\R{所}{ところ}にいて、1\R{分}{ぷん}に\qty{90}{m}ずつ\R{家}{いえ}に\R{近}{ちか}づくので、\R{位置}{い|ち}は $x$ \R{分}{ふん}で $90x$ だけ\R{小}{ちい}さくなります。$y=1500-90x$、\R{並}{なら}べかえて $y=-90x+1500$ です。\R{傾}{かたむ}き $-90$ は「1\R{分}{ぷん}に\qty{90}{m}、\R{家}{いえ}のほうへ\R{進}{すす}む」\R{速}{はや}さ、\R{切片}{せっ|ぺん}1500は\R{出発}{しゅっ|ぱつ}の\R{位置}{い|ち}です（3-5）。

\noindent\begin{minipage}[c]{0.55\linewidth}\centering
\begin{tikzpicture}
\begin{axis}[bkplot, width=62mm, height=58mm, xmin=0, xmax=17.5, ymin=0, ymax=1700, xtick={2,4,...,16}, ytick={300,600,...,1500}, xlabel={$x$（\R{分}{ふん}）}, ylabel={$y$（m）}]
  \addplot[bkcurve, domain=0:16.5] {60*x};
  \addplot[bkcurve2, domain=0:16.5] {-90*x+1500};
  \draw[bkGray, line width=0.6pt, densely dashed] (axis cs:10,0) -- (axis cs:10,600) -- (axis cs:0,600);
  \addplot[only marks, mark=*, mark size=2.2pt, bkRed] coordinates {(10,600)};
  \addplot[only marks, mark=*, mark size=1.6pt, bkInk] coordinates {(4,240)};
  \node[anchor=south west, font=\sffamily\footnotesize, text=bkRed!80!black, inner sep=1.5pt] at (axis cs:10.3,640) {$(10,\ 600)$};
  \node[anchor=north west, font=\sffamily\scriptsize, inner sep=1pt] at (axis cs:4.2,230) {$(4,\ 240)$};
  \node[anchor=west, font=\sffamily\footnotesize, text=bkBlue!80!black, inner sep=1pt] at (axis cs:15,1050) {ミナ};
  \node[anchor=west, font=\sffamily\footnotesize, text=bkOrange!80!black, inner sep=1pt] at (axis cs:4.6,1240) {ソラ};
\end{axis}
\end{tikzpicture}
\end{minipage}\hfill
\begin{minipage}[c]{0.42\linewidth}\raggedright
\R{青}{あお}がミナさんの $y=60x$、オレンジがソラさんの $y=-90x+1500$ です。\par\vspace{2pt}
2\R{本}{ほん}の\R{直線}{ちょく|せん}は、\R{点}{てん} $(10,\ 600)$ で\R{交}{まじ}わっています。10\R{分後}{ぷん|ご}、\R{家}{いえ}から\qty{600}{m}の\R{所}{ところ}で\R{出会}{で|あ}う、ということです。
\end{minipage}

\par\vspace{4pt}
\textbf{\R{線}{せん}の\R{上}{うえ}の\R{点}{てん}は、その\R{式}{しき}をみたす}\quad \R{青}{あお}い\R{線}{せん}の\R{上}{うえ}の\R{点}{てん} $(4,\ 240)$ を、ミナさんの\R{式}{しき}に\R{入}{い}れると $60\times4=240$ で、\R{等式}{とう|しき}が\R{成}{な}り\R{立}{た}ちます。「\R{線}{せん}の\R{上}{うえ}にある」ことと「その\R{式}{しき}をみたす」ことは、\R{同}{おな}じことです。でも、$(4,\ 240)$ をソラさんの\R{式}{しき}に\R{入}{い}れると、$-90\times4+1500=1140$ で、240になりません。この\R{点}{てん}は、オレンジの\R{線}{せん}の\R{上}{うえ}にはないのです（4\R{分後}{ふん|ご}、ソラさんは\qty{1140}{m}の\R{所}{ところ}）。

\textbf{\R{交}{まじ}わる\R{点}{てん}は、\R{両方}{りょう|ほう}をみたす}\quad $(10,\ 600)$ を\R{入}{い}れると、
\bkthrgline{ミナ：$60\times10=600$\qquad ソラ：$-90\times10+1500=-900+1500=600$}
\R{両方}{りょう|ほう}の\R{式}{しき}が\R{成}{な}り\R{立}{た}ちます。\R{交}{まじ}わる\R{点}{てん}は、\R{両方}{りょう|ほう}の\R{線}{せん}の\R{上}{うえ}にあるので、\R{両方}{りょう|ほう}の\R{式}{しき}を\R{同時}{どう|じ}にみたす\R{点}{てん}です。\R{交}{まじ}わる\R{点}{てん}を、\textbf{\R{交点}{こう|てん}}といいます。\bkidx{こうてん}{\R{交点}{こう|てん}}

\noindent\begin{minipage}[c]{0.46\linewidth}
このように、2つの\R{式}{しき}を\R{組}{くみ}にして、\R{両方}{りょう|ほう}を\R{同時}{どう|じ}に\R{成}{な}り\R{立}{た}たせる $x$ と $y$ をさがすとき、その\R{組}{くみ}の\R{式}{しき}を\textbf{\R{連立方程式}{れん|りつ|ほう|てい|しき}}といいます。\bkidx{れんりつほうていしき}{\R{連立方程式}{れん|りつ|ほう|てい|しき}}「\R{連立}{れん|りつ}」は、ならんで\R{立}{た}つという\R{意味}{い|み}です。\R{右}{みぎ}のカードのように、\R{左}{ひだり}に\R{大}{おお}きな\R{波}{なみ}かっこをつけて\R{書}{か}きます。
\end{minipage}\hfill
\begin{minipage}[c]{67mm}\bkkigocard{$\begin{cases}\,\cdots\\\,\cdots\end{cases}$}{\R{連立}{れん|りつ}の\R{波}{なみ}かっこ}{2つ（いくつか）の\R{式}{しき}を、\R{同時}{どう|じ}に\R{成}{な}り\R{立}{た}たせる\R{組}{くみ}だという\R{印}{しるし}。}{system of equations}{3-6}\end{minipage}

\[ \begin{cases} y=60x\\ y=-90x+1500 \end{cases} \]
2つの\R{式}{しき}を\R{同時}{どう|じ}に\R{成}{な}り\R{立}{た}たせる $x$ と $y$ の\R{組}{くみ}を、\R{連立方程式}{れん|りつ|ほう|てい|しき}の\textbf{\R{解}{かい}}といい、「$x=10$、$y=600$」や「$(x,\ y)=(10,\ 600)$」と\R{書}{か}きます。$x$ だけ、$y$ だけでは\R{解}{かい}になりません。2つで1\R{組}{くみ}です。

\bkhatenaref の\R{答}{こた}え：2\R{人}{にん}のグラフの\R{交点}{こう|てん}が、\R{出会}{で|あ}う\R{時間}{じ|かん}と\R{場所}{ば|しょ}です。10\R{分後}{ぷん|ご}、\R{家}{いえ}から\qty{600}{m}の\R{所}{ところ}で\R{出会}{で|あ}います。

\bksubsection{2}{\R{代入法}{だい|にゅう|ほう}：$y$ のかわりに\R{式}{しき}を\R{入}{い}れる}

\leavevmode\label{bk:3-6:dainyuu}グラフは、\R{交点}{こう|てん}の\R{目}{め}もりがちょうど\R{読}{よ}めるときには\R{便利}{べん|り}です。でも、\R{交点}{こう|てん}が\R{目}{め}もりのあいだにあると、\R{正確}{せい|かく}には\R{読}{よ}めません。そこで、\R{計算}{けい|さん}で\R{解}{と}く\R{方法}{ほう|ほう}を2つ\R{学}{まな}びます。

\bkthrgstep{1}{\R{出会}{で|あ}う\R{所}{ところ}では、2\R{人}{にん}の $y$ が\R{等}{ひと}しい}
\R{交点}{こう|てん}では、ミナさんの $y$ とソラさんの $y$ が\R{同}{おな}じ\R{数}{かず}です。ミナさんの\R{式}{しき}から、$y$ は $60x$ と\R{等}{ひと}しいので、ソラさんの\R{式}{しき}の $y$ のかわりに $60x$ を\R{入}{い}れます。
\begin{align*}
60x&=-90x+1500 && \text{\small（ソラさんの\R{式}{しき}の $y$ に $60x$ を\R{代入}{だい|にゅう}）}\\
60x+90x&=1500 && \text{\small（$-90x$ を\R{移項}{い|こう}。2-14）}\\
150x&=1500\\
x&=10 && \text{\small（\R{両辺}{りょう|へん}を150でわる）}
\end{align*}
$y$ が\R{消}{き}えて、$x$ だけの1\R{次方程式}{じ|ほう|てい|しき}になりました。これは2-16と\R{同}{おな}じ\R{形}{かたち}です。$x=10$ をミナさんの\R{式}{しき}に\R{入}{い}れて、$y=60\times10=600$。グラフで\R{読}{よ}んだ\R{答}{こた}えと\R{合}{あ}っています。

$150x$ の150は、2\R{人}{にん}が1\R{分}{ぷん}に\R{近}{ちか}づく\R{距離}{きょ|り} $60+90=150$（m）です。\qty{1500}{m}のあいだが、1\R{分}{ぷん}に\qty{150}{m}ずつちぢまるので、$1500\div150=10$ \R{分}{ぷん}で0になる、とも\R{考}{かんが}えられます（2-16の「\R{差}{さ}に\R{目}{め}をつける」）。

\bkthrgstep{2}{かっこごと\R{入}{い}れる}
\R{次}{つぎ}の\R{連立方程式}{れん|りつ|ほう|てい|しき}では、2つめの\R{式}{しき}が「$y=$」の\R{形}{かたち}になっていません。
\[ \begin{cases} y=2x-1 & \cdots\text{①}\\ 3x+2y=12 & \cdots\text{②} \end{cases} \]
①から、$y$ は $2x-1$ と\R{等}{ひと}しいので、②の $y$ のかわりに $2x-1$ を\R{入}{い}れます。②の $2y$ は「$2\times y$」なので、$y$ の\R{箱}{はこ}に $2x-1$ をまるごと\R{入}{い}れて、$2\times(2x-1)$ にします（2-13）。

\begin{bkfigure}
\begin{tikzpicture}[x=1mm, y=1mm]
  \node[font=\Large, inner sep=1pt, anchor=east] (a) at (0,0) {$3x+2\,$};
  \node[draw=bkOrange!80!black, fill=bkOrange!12, line width=0.9pt, rounded corners=1mm, minimum width=11mm, minimum height=8mm, font=\Large, anchor=west, inner sep=1pt] (b) at (a.east) {$y$};
  \node[font=\Large, anchor=west, inner sep=1pt] (c) at ([xshift=1mm]b.east) {$=12$};
  \draw[bkGray, line width=0.9pt, -{Stealth[length=5pt]}] ([xshift=4mm]c.east) -- ([xshift=16mm]c.east);
  \node[font=\Large, inner sep=1pt, anchor=west] (d) at ([xshift=20mm]c.east) {$3x+2\,$};
  \node[draw=bkOrange!80!black, fill=bkOrange!12, line width=0.9pt, rounded corners=1mm, minimum height=8mm, font=\Large, anchor=west, inner sep=1.5pt] (e) at (d.east) {$(2x-1)$};
  \node[font=\Large, anchor=west, inner sep=1pt] at ([xshift=1mm]e.east) {$=12$};
  \node[font=\sffamily\footnotesize, text=bkOrange!80!black, anchor=north] at ([yshift=-1.5mm]b.south) {$y$ の\R{箱}{はこ}};
  \node[font=\sffamily\footnotesize, text=bkOrange!80!black, anchor=north] at ([yshift=-1.5mm]e.south) {かっこごと\R{入}{い}れる};
\end{tikzpicture}
\bkfigcaption{$y$ の\R{箱}{はこ}に、$2x-1$ をまるごと\R{入}{い}れる。}
\end{bkfigure}

\begin{align*}
3x+2(2x-1)&=12 && \text{\small（②の $y$ に①を\R{代入}{だい|にゅう}）}\\
3x+4x-2&=12 && \text{\small（かっこをはずす。$2\times2x=4x$、$2\times(-1)=-2$）}\\
7x&=14 && \text{\small（$3x+4x=7x$。$-2$ を\R{移項}{い|こう}して $12+2=14$）}\\
x&=2
\end{align*}
$x=2$ を①に\R{入}{い}れると、$y=2\times2-1=3$。\R{解}{かい}は $x=2$、$y=3$ です。

\textbf{\R{確}{たし}かめ}\quad ②に\R{入}{い}れると、$3\times2+2\times3=6+6=12$。①も②も\R{成}{な}り\R{立}{た}つので、\R{正}{ただ}しい\R{解}{かい}です。\R{連立方程式}{れん|りつ|ほう|てい|しき}の\R{確}{たし}かめは、\R{必}{かなら}ず\R{両方}{りょう|ほう}の\R{式}{しき}でします。

\noindent このように、1つの\R{式}{しき}を「$y=\cdots$」（または「$x=\cdots$」）の\R{形}{かたち}にして、もう1つの\R{式}{しき}に\R{代入}{だい|にゅう}し、\R{文字}{も|じ}を1つ\R{消}{け}して\R{解}{と}く\R{方法}{ほう|ほう}を、\textbf{\R{代入法}{だい|にゅう|ほう}}といいます。\bkidx{だいにゅうほう}{\R{代入法}{だい|にゅう|ほう}}

\begin{bkcenter}
\begin{bktable}{colspec={Q[c,m,wd=10mm] Q[l,m,wd=104mm]}, rowsep=1.5pt, cells={font=\small}}
① & 1つの\R{式}{しき}を「$y=\cdots$」の\R{形}{かたち}にする（はじめからその\R{形}{かたち}なら、そのまま）。 \\
② & もう1つの\R{式}{しき}の $y$ に、かっこごと\R{代入}{だい|にゅう}する。$x$ だけの1\R{次方程式}{じ|ほう|てい|しき}になる。 \\
③ & 1\R{次方程式}{じ|ほう|てい|しき}を\R{解}{と}いて $x$ を\R{求}{もと}める（2-16）。 \\
④ & $x$ を「$y=\cdots$」の\R{式}{しき}に\R{入}{い}れて $y$ を\R{求}{もと}め、\R{両方}{りょう|ほう}の\R{式}{しき}で\R{確}{たし}かめる。 \\
\end{bktable}
\bktabcaption{\R{代入法}{だい|にゅう|ほう}の\R{手順}{て|じゅん}}
\end{bkcenter}

\bksubsection{3}{\R{加減法}{か|げん|ほう}：\R{式}{しき}どうしをたす・ひく}

\noindent\begin{minipage}[c]{0.44\linewidth}\label{bk:3-6:kagen}
くだもの\R{屋}{や}さんで、りんごとみかんを\R{買}{か}いました。\R{右}{みぎ}の2つのトレーの\R{値段}{ね|だん}から、りんご1\R{個}{こ}とみかん1\R{個}{こ}の\R{値段}{ね|だん}がわかるでしょうか。
\end{minipage}\hfill
\begin{minipage}[c]{0.54\linewidth}\centering
\begin{tikzpicture}[x=1mm, y=1mm]
  \bkthrgtray{(3,16)}{2}{3}{420}
  \bkthrgtray{(3,0)}{2}{1}{300}
  \node[font=\sffamily\scriptsize, anchor=east] at (1.5,22) {A};
  \node[font=\sffamily\scriptsize, anchor=east] at (1.5,6) {B};
\end{tikzpicture}
\end{minipage}

\par\vspace{4pt}
\bkthrgstep{1}{くらべて\R{考}{かんが}える}
AとBをくらべると、りんごの\R{数}{かず}は\R{同}{おな}じ2\R{個}{こ}で、ちがうのはみかんの\R{数}{かず}だけです。Aのほうがみかんが $3-1=2$ \R{個}{こ}\R{多}{おお}く、\R{値段}{ね|だん}は $420-300=120$ \R{円}{えん}\R{高}{たか}いので、みかん2\R{個}{こ}で120\R{円}{えん}。みかん1\R{個}{こ}は $120\div2=60$ \R{円}{えん}です。Bから、りんご2\R{個}{こ}は $300-60=240$ \R{円}{えん}なので、りんご1\R{個}{こ}は $240\div2=120$ \R{円}{えん}です。

\bkthrgstep{2}{\R{式}{しき}で\R{考}{かんが}える}
りんご1\R{個}{こ}を $x$ \R{円}{えん}、みかん1\R{個}{こ}を $y$ \R{円}{えん}とすると、
\[ \begin{cases} 2x+3y=420 & \cdots\text{①}\\ 2x+y=300 & \cdots\text{②} \end{cases} \]
①から②を、\R{左辺}{さ|へん}どうし、\R{右辺}{う|へん}どうしでひきます。\R{同}{おな}じ\R{文字}{も|じ}の\R{項}{こう}をたてにそろえて\R{書}{か}くと、ひき\R{算}{ざん}の\R{筆算}{ひっ|さん}のようにできます。
\[ \begin{array}{r@{\;}r@{\;}c@{\;}r@{\;}c@{\;}r}
 & 2x & + & 3y & = & 420\\
-) & 2x & + & y & = & 300\\ \hline
 & & & 2y & = & 120
\end{array} \]
$2x-2x=0$ で $x$ が\R{消}{き}え、$3y-y=2y$、$420-300=120$。$2y=120$ から $y=60$。これを②に\R{入}{い}れて $2x+60=300$、$2x=240$、$x=120$ です。\R{確}{たし}かめ：①は $2\times120+3\times60=240+180=420$ で\R{成}{な}り\R{立}{た}ちます。

\bkthrgstep{3}{なぜ\R{式}{しき}どうしをひいてよいの？}
\noindent\begin{minipage}[c]{0.52\linewidth}\raggedright
\R{等式}{とう|しき}は、つり\R{合}{あ}った\R{天}{てん}びんでした（2-14）。①の\R{天}{てん}びんは、\R{左}{ひだり}に $2x+3y$、\R{右}{みぎ}に420\R{円}{えん}がのってつり\R{合}{あ}っています。②から、$2x+y$ と300\R{円}{えん}は\R{同}{おな}じ\R{重}{おも}さです。①の\R{左}{ひだり}の\R{皿}{さら}から $2x+y$ を、\R{右}{みぎ}の\R{皿}{さら}から300\R{円}{えん}をとると、\R{同}{おな}じ\R{重}{おも}さをとったので、つり\R{合}{あ}ったままです。\R{残}{のこ}りは、\R{左}{ひだり}が $2y$、\R{右}{みぎ}が120\R{円}{えん}。だから $2y=120$ です。
\end{minipage}\hfill
\begin{minipage}[c]{0.45\linewidth}\centering
\begin{tikzpicture}[x=1mm, y=1mm]
  \fill[bkGray!70, rounded corners=0.6mm] (-10,0) rectangle (10,3);
  \fill[bkGray!80] (-1.4,3) rectangle (1.4,17);
  \fill[bkInk!75, rounded corners=0.8mm] (-22,17) rectangle (22,19);
  \fill[bkRed] (0,18) circle (1.3);
  \foreach \s in {-1,1} {\fill[bkGray!80] ({\s*18-0.8},19) rectangle ({\s*18+0.8},21.5);
    \fill[bkGray!55, rounded corners=0.5mm] ({\s*18-12.5},21.5) rectangle ({\s*18+12.5},23);}
  % 左の皿：りんご2こ・みかん3こ（点線の中は 2x+y）
  \bkapple{(-28,26)}{2.3}{bkAppleRed}\bkapple{(-23,26)}{2.3}{bkAppleRed}
  \bkthrgmikan{(-18,25.6)}{2.2}
  \bkthrgmikan{(-12.6,25.6)}{2.2}\bkthrgmikan{(-8,25.6)}{2.2}
  \draw[bkBlue, line width=0.8pt, rounded corners=1mm, densely dashed] (-31,22.4) rectangle (-15.3,29.8);
  % 右の皿：300円と120円（点線の中は300円）
  \foreach \x/\t in {12.2/300, 21.6/120} {\fill[bkYellow!60] (\x-4.4,23) rectangle (\x+4.4,28.4);
    \draw[bkYellow!60!black, line width=0.4pt] (\x-4.4,23) rectangle (\x+4.4,28.4);
    \node[font=\sffamily\scriptsize, inner sep=0pt] at (\x,25.7) {\t};}
  \draw[bkBlue, line width=0.8pt, rounded corners=1mm, densely dashed] (7.2,22.4) rectangle (17.2,29.4);
  \node[font=\sffamily\scriptsize, anchor=south] at (17,30.3) {\R{合}{あ}わせて420\R{円}{えん}};
  \node[font=\sffamily\scriptsize, text=bkBlue!80!black, anchor=north, align=center] at (0,-1) {\R{点線}{てん|せん}の\R{中}{なか}は\R{同}{おな}じ\R{重}{おも}さ\\（②：$2x+y$ と300\R{円}{えん}）};
\end{tikzpicture}
\end{minipage}

\par\vspace{4pt}
\R{図}{ず}では、\R{右}{みぎ}の\R{皿}{さら}の420\R{円}{えん}を、300\R{円}{えん}と120\R{円}{えん}に\R{分}{わ}けてかいています。\R{点線}{てん|せん}の\R{中}{なか}どうしをとると、\R{左}{ひだり}にはみかん2\R{個}{こ}、\R{右}{みぎ}には120\R{円}{えん}が\R{残}{のこ}ります。\R{等}{ひと}しいものから\R{等}{ひと}しいものをひいても、\R{残}{のこ}りは\R{等}{ひと}しいのです。たすときも\R{同}{おな}じで、\R{等}{ひと}しいものに\R{等}{ひと}しいものをたしても、つり\R{合}{あ}ったままです。

\bkthrgstep{4}{たして\R{消}{け}す}
「2つの\R{数}{かず}の\R{和}{わ}が7、\R{差}{さ}が1」のとき、\R{大}{おお}きいほうを $x$、\R{小}{ちい}さいほうを $y$ とすると、
\[ \begin{cases} x+y=7 & \cdots\text{①}\\ x-y=1 & \cdots\text{②} \end{cases}
\qquad\qquad
\begin{array}{r@{\;}r@{\;}c@{\;}r@{\;}c@{\;}r}
 & x & + & y & = & 7\\
+) & x & - & y & = & 1\\ \hline
 & 2x & & & = & 8
\end{array} \]
$+y$ と $-y$ は、たすと0になって\R{消}{き}えます。$2x=8$ から $x=4$。①に\R{入}{い}れて $4+y=7$、$y=3$。\R{確}{たし}かめ：②は $4-3=1$。$y$ の\R{符号}{ふ|ごう}が\R{同}{おな}じならひき、\R{反対}{はん|たい}ならたすと、$y$ が\R{消}{き}えます。

\bkthrgstep{5}{\R{係数}{けい|すう}がそろっていないときは、かけてそろえる}
\R{次}{つぎ}の\R{連立方程式}{れん|りつ|ほう|てい|しき}を\R{解}{と}きます。
\[ \begin{cases} x+2y=8 & \cdots\text{①}\\ 3x-y=3 & \cdots\text{②} \end{cases} \]
このままたしても、ひいても、どちらの\R{文字}{も|じ}も\R{消}{き}えません。そこで、②の\R{両辺}{りょう|へん}に2をかけて、$y$ の\R{係数}{けい|すう}を $-2$ にします。\R{両辺}{りょう|へん}に\R{同}{おな}じ\R{数}{かず}をかけても、\R{等式}{とう|しき}は\R{成}{な}り\R{立}{た}ちます（2-14）。\R{全部}{ぜん|ぶ}の\R{項}{こう}にかけるので、\R{右辺}{う|へん}の3にもかけます。
\[ \text{②}\times2\qquad 6x-2y=6\quad\cdots\text{②}' \qquad\qquad
\begin{array}{r@{\;}r@{\;}c@{\;}r@{\;}c@{\;}r}
 & x & + & 2y & = & 8\\
+) & 6x & - & 2y & = & 6\\ \hline
 & 7x & & & = & 14
\end{array} \]
$x=2$。①に\R{入}{い}れて $2+2y=8$、$2y=6$、$y=3$。\R{確}{たし}かめ：②は $3\times2-3=3$ で\R{成}{な}り\R{立}{た}ちます。②$'$ は「②を2\R{倍}{ばい}した\R{式}{しき}」の\R{印}{しるし}です（「ダッシュ」と\R{読}{よ}みます）。

\noindent このように、\R{式}{しき}どうしをたしたりひいたりして\R{文字}{も|じ}を1つ\R{消}{け}す\R{方法}{ほう|ほう}を、\textbf{\R{加減法}{か|げん|ほう}}といいます。\bkidx{かげんほう}{\R{加減法}{か|げん|ほう}}「\R{加}{か}」はたし\R{算}{ざん}、「\R{減}{げん}」はひき\R{算}{ざん}です。

\textbf{\R{代入法}{だい|にゅう|ほう}で\R{解}{と}いても\R{同}{おな}じ}\quad くだもの\R{屋}{や}さんの\R{問題}{もん|だい}を\R{代入法}{だい|にゅう|ほう}で\R{解}{と}いてみます。②の $2x$ を\R{移項}{い|こう}して $y=300-2x$。①に\R{代入}{だい|にゅう}すると、
\[ 2x+3(300-2x)=420 \]
かっこをはずして $2x+900-6x=420$、まとめて $-4x=-480$、\R{両辺}{りょう|へん}を $-4$ でわって $x=120$。
$y=300-2\times120=60$ で、\R{同}{おな}じ\R{答}{こた}えです。どちらの\R{方法}{ほう|ほう}でも\R{解}{と}けます。「$y=\cdots$」の\R{形}{かたち}の\R{式}{しき}があれば\R{代入法}{だい|にゅう|ほう}、\R{係数}{けい|すう}がそろっていたり、そろえやすかったりすれば\R{加減法}{か|げん|ほう}が\R{楽}{らく}です。

\bksubsection{4}{\R{平行}{へい|こう}なら\R{交}{まじ}わらない、\R{重}{かさ}なれば\R{無数}{む|すう}}

\leavevmode\label{bk:3-6:heikou}2-16の「もっと\R{知}{し}りたい」で、\R{兄}{あに}も\R{弟}{おとうと}も\R{秒速}{びょう|そく}\qty{2}{m}だったら、\R{方程式}{ほう|てい|しき} $2t=10+2t$ には\R{解}{かい}がない、という\R{話}{はなし}をしました。「グラフでは、2\R{本}{ほん}の\R{線}{せん}が\R{同}{おな}じかたむきで\R{平行}{へい|こう}になり、いつまでも\R{交}{まじ}わりません（3-6）」と\R{書}{か}いてありました。ここで、そのわけを\R{確}{たし}かめます。

\noindent\begin{minipage}[c]{0.5\linewidth}\centering
\begin{tikzpicture}
\begin{axis}[bkplot, width=60mm, height=52mm, xmin=0, xmax=8.6, ymin=0, ymax=29, xtick={1,...,8}, ytick={5,10,...,25}, xlabel={$x$（s）}, ylabel={$y$（m）}]
  \addplot[bkcurve, domain=0:8.3] {2*x};
  \addplot[bkcurve2, domain=0:8.3] {2*x+10};
  \foreach \xx in {2,6} {\edef\tmp{\noexpand\draw[{Stealth[length=3pt]}-{Stealth[length=3pt]}, bkRed, line width=0.6pt] (axis cs:\xx,{2*\xx+0.5}) -- (axis cs:\xx,{2*\xx+9.5});}\tmp}
  \node[anchor=west, font=\sffamily\scriptsize, text=bkRed!80!black, inner sep=1pt] at (axis cs:6.1,17) {\qty{10}{m}};
  \node[anchor=west, font=\sffamily\footnotesize, text=bkBlue!80!black, inner sep=1pt] at (axis cs:7,11.5) {\R{兄}{あに}};
  \node[anchor=east, font=\sffamily\footnotesize, text=bkOrange!80!black, inner sep=1pt] at (axis cs:4,22) {\R{弟}{おとうと}};
\end{axis}
\end{tikzpicture}
\end{minipage}\hfill
\begin{minipage}[c]{0.47\linewidth}
$x$ \R{秒後}{びょう|ご}の\R{兄}{あに}の\R{位置}{い|ち}を $y=2x$、\R{弟}{おとうと}の\R{位置}{い|ち}を $y=2x+10$ とします。\R{傾}{かたむ}きはどちらも2で\R{同}{おな}じ、\R{切片}{せっ|ぺん}は0と10でちがいます。\par\vspace{2pt}
2\R{本}{ほん}の\R{線}{せん}は、どこでもたてに\qty{10}{m}はなれたまま\R{進}{すす}みます。\R{同}{おな}じ\R{傾}{かたむ}きなので、\R{近}{ちか}づきも、はなれもしないからです。
\end{minipage}

\par\vspace{4pt}
\R{代入法}{だい|にゅう|ほう}で\R{解}{と}こうとすると、
\[ 2x=2x+10\ \to\ 2x-2x=10\ \to\ 0=10 \]
$x$ が\R{消}{き}えて、\R{成}{な}り\R{立}{た}たない\R{式}{しき} $0=10$ が\R{残}{のこ}ります。$x$ にどんな\R{数}{かず}を\R{入}{い}れても、0は10になりません。だから、この\R{連立方程式}{れん|りつ|ほう|てい|しき}には\textbf{\R{解}{かい}がありません}。グラフの\R{平行}{へい|こう}な2\R{本}{ほん}の\R{線}{せん}には、\R{交点}{こう|てん}がありません。\R{式}{しき}で\R{見}{み}ても、グラフで\R{見}{み}ても、\R{同}{おな}じ\R{結論}{けつ|ろん}です。

\textbf{\R{重}{かさ}なるとき}\quad こんどは、\R{次}{つぎ}の\R{連立方程式}{れん|りつ|ほう|てい|しき}を\R{代入法}{だい|にゅう|ほう}で\R{解}{と}いてみます。
\[ \begin{cases} y=2x+1 & \cdots\text{①}\\ 4x-2y=-2 & \cdots\text{②} \end{cases} \]
②に①を\R{代入}{だい|にゅう}すると、
\[ 4x-2(2x+1)=-2\ \to\ 4x-4x-2=-2\ \to\ -2=-2 \]
となります。$x$ が\R{消}{き}えて、いつでも\R{成}{な}り\R{立}{た}つ\R{式}{しき}が\R{残}{のこ}りました。$x$ がどんな\R{数}{かず}でも、$y=2x+1$ なら②も\R{成}{な}り\R{立}{た}つ、ということです。②を\R{変形}{へん|けい}すると、$-2y=-4x-2$、\R{両辺}{りょう|へん}を $-2$ でわって $y=2x+1$。②は、①と\R{同}{おな}じ\R{直線}{ちょく|せん}だったのです。2\R{本}{ほん}の\R{線}{せん}がぴったり\R{重}{かさ}なるので、\R{直線}{ちょく|せん}の\R{上}{うえ}のどの\R{点}{てん}も\R{交点}{こう|てん}です。\R{解}{かい}は、$(0,\ 1)$、$(1,\ 3)$、$(2,\ 5)$、……と\textbf{\R{無数}{む|すう}}（\R{数}{かぞ}えきれないほどたくさん）にあります。

\begin{bkfigure}
\begin{tikzpicture}
\begin{axis}[bkplot, name=h1, width=44mm, height=42mm, xmin=-0.5, xmax=4.4, ymin=-0.5, ymax=5.4, xtick={1,2,3,4}, ytick={1,...,5}, xlabel={$x$}, ylabel={$y$}, title={\sffamily\small 1\R{点}{てん}で\R{交}{まじ}わる}, title style={yshift=3mm}]
  \addplot[bkcurve, domain=0:2.2] {2*x+1};
  \addplot[bkcurve2, domain=0:4.2] {-x+4};
  \addplot[only marks, mark=*, mark size=2pt, bkRed] coordinates {(1,3)};
\end{axis}
\begin{axis}[bkplot, name=h2, at={($(h1.east)+(9mm,0)$)}, anchor=west, width=44mm, height=42mm, xmin=-0.5, xmax=4.4, ymin=-0.5, ymax=5.4, xtick={1,2,3,4}, ytick={1,...,5}, xlabel={$x$}, ylabel={$y$}, title={\sffamily\small \R{平行}{へい|こう}}, title style={yshift=3mm}]
  \addplot[bkcurve, domain=0:2.2] {2*x+1};
  \addplot[bkcurve2, domain=1.5:4.2] {2*x-3};
\end{axis}
\begin{axis}[bkplot, name=h3, at={($(h2.east)+(9mm,0)$)}, anchor=west, width=44mm, height=42mm, xmin=-0.5, xmax=4.4, ymin=-0.5, ymax=5.4, xtick={1,2,3,4}, ytick={1,...,5}, xlabel={$x$}, ylabel={$y$}, title={\sffamily\small \R{重}{かさ}なる}, title style={yshift=3mm}]
  \addplot[bkcurve2, line width=3.4pt, domain=0:2.2] {2*x+1};
  \addplot[bkcurve, line width=1.2pt, domain=0:2.2] {2*x+1};
\end{axis}
\end{tikzpicture}
\bkfigcaption{\R{左}{ひだり}：$y=2x+1$ と $y=-x+4$。\R{中}{なか}：$y=2x+1$ と $y=2x-3$。\R{右}{みぎ}：$y=2x+1$ と $4x-2y=-2$（\R{同}{おな}じ\R{直線}{ちょく|せん}）。}
\end{bkfigure}

\begin{bkcenter}
\begin{bktable}{colspec={Q[l,m,wd=42mm] Q[c,m,wd=26mm] Q[l,m,wd=46mm]}, rowsep=1.5pt, cells={font=\small}, row{1}={font=\sffamily\bfseries\small}}
2\R{本}{ほん}の\R{直線}{ちょく|せん} & グラフ & \R{連立方程式}{れん|りつ|ほう|てい|しき}の\R{解}{かい} \\
\R{傾}{かたむ}きがちがう & 1\R{点}{てん}で\R{交}{まじ}わる & 1\R{組}{くみ}だけ \\
\R{傾}{かたむ}きが\R{同}{おな}じで、\R{切片}{せっ|ぺん}がちがう & \R{平行}{へい|こう} & ない（$0=10$ のような\R{式}{しき}が\R{残}{のこ}る） \\
\R{傾}{かたむ}きも\R{切片}{せっ|ぺん}も\R{同}{おな}じ & \R{重}{かさ}なる & \R{無数}{む|すう}（$-2=-2$ のような\R{式}{しき}が\R{残}{のこ}る） \\
\end{bktable}
\bktabcaption{2つの\R{式}{しき}を「$y=ax+b$」の\R{形}{かたち}にすると、\R{解}{かい}がいくつあるか、\R{解}{と}く\R{前}{まえ}にわかる。}
\end{bkcenter}

\bksubsection{5}{\R{表計算}{ひょう|けい|さん}で\R{交点}{こう|てん}をさがす}

\leavevmode\label{bk:3-6:hyou}ミナさんが\R{分速}{ふん|そく}\qty{70}{m}で\R{歩}{ある}いたらどうなるでしょう。\R{連立方程式}{れん|りつ|ほう|てい|しき}は $y=70x$ と $y=-90x+1500$ です。2-16のように、\R{表}{ひょう}でためして\R{交点}{こう|てん}に\R{近}{ちか}づいてみます。\R{表計算}{ひょう|けい|さん}（2-18）を\R{使}{つか}うと\R{楽}{らく}です。D\R{列}{れつ}には、2\R{人}{にん}のあいだの\R{距離}{きょ|り}（ソラ－ミナ）を\R{計算}{けい|さん}させます。

\begin{bkfigure}
\begin{tikzpicture}[x=1mm, y=1mm]
  \node[anchor=north west, inner sep=0pt] at (0,0) {\bksheet{4}{13}{%
1 & $x$（\R{分}{ふん}） & ミナ & ソラ & \R{差}{さ} \\
2 & 8 & 560 & 780 & 220 \\
3 & 9 & 630 & 690 & 60 \\
4 & 10 & 700 & 600 & $-100$ \\
5 & 11 & 770 & 510 & $-260$ \\}};
  \node[anchor=north west, align=left, font=\sffamily\small, inner sep=0pt] at (82,-6) {B2：\texttt{=70*A2}\\[1.5mm]C2：\texttt{=1500-90*A2}\\[1.5mm]D2：\texttt{=C2-B2}\\[1.5mm]を\R{入}{い}れて、\R{下}{した}へコピー};
\end{tikzpicture}
\bkfigcaption{\R{差}{さ}が\R{正}{せい}から\R{負}{ふ}に\R{変}{か}わる9\R{分}{ふん}と10\R{分}{ぷん}のあいだに、\R{交点}{こう|てん}がある。}
\end{bkfigure}

\R{差}{さ}がプラスのうちは、ソラさんがまだ\R{家}{いえ}から\R{遠}{とお}い\R{所}{ところ}にいて、2\R{人}{にん}は\R{出会}{で|あ}っていません。マイナスになったら、すれちがった\R{後}{あと}です。だから、\R{差}{さ}が0になる\R{交点}{こう|てん}は、9\R{分}{ふん}と10\R{分}{ぷん}のあいだにあります。A\R{列}{れつ}を9.0、9.1、9.2、……と0.1ずつにすると、9.3\R{分}{ぷん}で\R{差}{さ}が12、9.4\R{分}{ふん}で $-4$ になり、9.3\R{分}{ぷん}と9.4\R{分}{ふん}のあいだまでしぼれます。

\R{代入法}{だい|にゅう|ほう}で\R{解}{と}くと、
\[ 70x=-90x+1500,\qquad 160x=1500,\qquad x=1500\div160=9.375 \]
で、$y=70\times9.375=656.25$ です。9.375\R{分}{ふん}の0.375\R{分}{ふん}を\R{秒}{びょう}にすると、$\qty{0.375}{min}\times\qty{60}{s/min}=\qty{22.5}{s}$（2-10）。9\R{分}{ふん}22.5\R{秒}{びょう}です。\R{表計算}{ひょう|けい|さん}は\R{答}{こた}えのありかを\R{見}{み}つけるのが\R{得意}{とく|い}で、\R{式}{しき}の\R{計算}{けい|さん}はぴったりの\R{答}{こた}えを\R{出}{だ}すのが\R{得意}{とく|い}です。\R{両方}{りょう|ほう}を\R{使}{つか}うと、まちがいにも\R{気}{き}づきやすくなります。

\bksubsection{6}{\R{係数}{けい|すう}を\R{表}{ひょう}にならべる：\R{行列}{ぎょう|れつ}の\R{種}{たね}}

\leavevmode\label{bk:3-6:keisuu}くだもの\R{屋}{や}さんの\R{連立方程式}{れん|りつ|ほう|てい|しき}をもう\R{一度}{いち|ど}\R{見}{み}ます。
\[ \begin{cases} 2x+3y=420\\ 2x+y=300 \end{cases} \]
\R{加減法}{か|げん|ほう}で\R{計算}{けい|さん}しているあいだ、$x$ や $y$ の\R{文字}{も|じ}は\R{動}{うご}きません。\R{変}{か}わるのは、\R{文字}{も|じ}にかかっている\R{数}{かず}（\textbf{\R{係数}{けい|すう}}）と\R{右辺}{う|へん}の\R{数}{かず}だけです。\bkidx{けいすう}{\R{係数}{けい|すう}}そこで、\R{文字}{も|じ}を\R{書}{か}かずに、\R{数}{かず}だけを\R{表}{ひょう}にならべてみます。1\R{行}{ぎょう}が1つの\R{式}{しき}です。たとえば1\R{行}{ぎょう}め「2、3、420」は、$2x+3y=420$ のことです。

\begin{bkfigure}
\begin{tikzpicture}[x=1mm, y=1mm]
  \foreach \xx in {0,48,96} {\node[font=\small] at (\xx+4.5,3) {$x$};\node[font=\small] at (\xx+13.5,3) {$y$};\node[font=\sffamily\scriptsize] at (\xx+22.5,3) {\R{右辺}{う|へん}};}
  \bkthrgtab{(0,0)}{2}{3}{420}{2}{1}{300}
  \bkthrgtab{(48,0)}{2}{3}{420}{0}{1}{60}
  \bkthrgtab{(96,0)}{1}{0}{120}{0}{1}{60}
  \draw[bkGray, line width=0.9pt, -{Stealth[length=4pt]}] (29,-6) -- (46,-6);
  \draw[bkGray, line width=0.9pt, -{Stealth[length=4pt]}] (77,-6) -- (94,-6);
  \node[font=\sffamily\scriptsize, align=center, anchor=north] at (37.5,-14.5) {2\R{行}{ぎょう}め－1\R{行}{ぎょう}め\\を $-2$ でわる};
  \node[font=\sffamily\scriptsize, align=center, anchor=north] at (85.5,-14.5) {1\R{行}{ぎょう}め－2\R{行}{ぎょう}め×3\\を2でわる};
  \node[font=\sffamily\scriptsize, anchor=north] at (13.5,-14.5) {はじめ};
  \node[font=\sffamily\scriptsize, anchor=north, align=center] at (109.5,-14.5) {$x=120$\\$y=60$};
\end{tikzpicture}
\bkfigcaption{\R{数}{かず}の\R{表}{ひょう}だけで\R{加減法}{か|げん|ほう}をする。\R{最後}{さい|ご}の\R{表}{ひょう}は「$1x+0y=120$」「$0x+1y=60$」と\R{読}{よ}む。}
\end{bkfigure}

1つめの\R{矢印}{や|じるし}：2\R{行}{ぎょう}めから1\R{行}{ぎょう}めをひくと、「$0$、$-2$、$-120$」（$2-2=0$、$1-3=-2$、$300-420=-120$）。これを $-2$ でわって「0、1、60」。$0x+1y=60$、つまり $y=60$ です。\par
2つめの\R{矢印}{や|じるし}：1\R{行}{ぎょう}めから、2\R{行}{ぎょう}めを3\R{倍}{ばい}した「0、3、180」をひくと「2、0、240」。2でわって「1、0、120」。$x=120$ です。

\R{表}{ひょう}の\R{行}{ぎょう}どうしをひいたり、\R{行}{ぎょう}に\R{数}{かず}をかけたりするのは、\R{加減法}{か|げん|ほう}そのものです。\R{文字}{も|じ}を\R{書}{か}かなくても\R{解}{と}けるので、コンピュータはこのやり\R{方}{かた}で、\R{文字}{も|じ}が\R{何千}{なん|ぜん}もある\R{連立方程式}{れん|りつ|ほう|てい|しき}を\R{解}{と}いています。この\R{数}{かず}の\R{表}{ひょう}を、\R{大学}{だい|がく}の\R{数学}{すう|がく}では\textbf{\R{行列}{ぎょう|れつ}}とよび、7-1で\R{学}{まな}びます。

\bksubsection{7}{\R{回路}{かい|ろ}の2つの\R{電流}{でん|りゅう}を、いっしょに\R{求}{もと}める}

\leavevmode\label{bk:3-6:kairo}2-16の\R{最後}{さい|ご}に「\R{電気}{でん|き}の\R{回路}{かい|ろ}では、わからない\R{電流}{でん|りゅう}をいくつも\R{文字}{も|じ}にして、\R{方程式}{ほう|てい|しき}を\R{何本}{なん|ぼん}もいっしょに\R{解}{と}きます（\R{連立方程式}{れん|りつ|ほう|てい|しき}、3-6）」と\R{書}{か}きました。\R{下}{した}の\R{回路}{かい|ろ}には、\R{電源}{でん|げん}（\R{電源装置}{でん|げん|そう|ち}）が2つあります。\R{左}{ひだり}の\R{電源}{でん|げん}が\R{流}{なが}す\R{電流}{でん|りゅう}を $I_1$、\R{右}{みぎ}の\R{電源}{でん|げん}が\R{流}{なが}す\R{電流}{でん|りゅう}を $I_2$（\R{単位}{たん|い}はA）とします。

\begin{bkfigure}
\begin{tikzpicture}[x=1mm, y=1mm]
  \bkthrgcircuit{\qty{9}{V}}{\qty{2}{\ohm}}{\qty{2}{\ohm}}{\qty{2}{\ohm}}{\qty{6}{V}}
  \draw[bkGreen!70!black, line width=0.8pt, densely dashed, -{Stealth[length=3.5pt]}] (27,24) arc (20:330:7.5);
  \node[font=\sffamily\scriptsize, text=bkGreen!50!black] at (20,21.4) {\R{左}{ひだり}の\R{輪}{わ}};
  \draw[bkGreen!70!black, line width=0.8pt, densely dashed, -{Stealth[length=3.5pt]}] (53,24) arc (160:-150:7.5);
  \node[font=\sffamily\scriptsize, text=bkGreen!50!black] at (60,21.4) {\R{右}{みぎ}の\R{輪}{わ}};
\end{tikzpicture}
\bkfigcaption{\R{電源}{でん|げん}の\R{長}{なが}い\R{線}{せん}が＋。\R{抵抗}{てい|こう}は\R{四角}{し|かく}でかく。\R{赤}{あか}い\R{矢印}{や|じるし}が\R{電流}{でん|りゅう}の\R{向}{む}き。}
\end{bkfigure}

\R{式}{しき}を\R{作}{つく}るために、\R{電流}{でん|りゅう}の3つのきまりを\R{使}{つか}います（\R{中}{ちゅう}2の\R{理科}{り|か}と2-15）。

\noindent\textbf{きまり1：\R{合}{あ}わさった\R{電流}{でん|りゅう}は、たし\R{算}{ざん}}\par
\R{上}{うえ}の\R{線}{せん}で、\R{左}{ひだり}から\R{来}{き}た $I_1$ と\R{右}{みぎ}から\R{来}{き}た $I_2$ が\R{合}{あ}わさって、まん\R{中}{なか}の\R{抵抗}{てい|こう}を\R{下}{した}へ\R{流}{なが}れます。\R{電流}{でん|りゅう}は\R{途中}{と|ちゅう}でなくならないので、まん\R{中}{なか}の\R{電流}{でん|りゅう}は $I_1+I_2$ です。

\noindent\textbf{きまり2：\R{抵抗}{てい|こう}にかかる\R{電圧}{でん|あつ}は、\R{抵抗}{てい|こう}×\R{電流}{でん|りゅう}}\par
オームの\R{法則}{ほう|そく} $V=RI$ です（2-15）。\R{左}{ひだり}の\R{抵抗}{てい|こう}には $2I_1$、まん\R{中}{なか}には $2(I_1+I_2)$、\R{右}{みぎ}には $2I_2$（V）の\R{電圧}{でん|あつ}がかかります。

\noindent\textbf{きまり3：1つの\R{輪}{わ}では、\R{電源}{でん|げん}の\R{電圧}{でん|あつ}＝\R{抵抗}{てい|こう}の\R{電圧}{でん|あつ}の\R{和}{わ}}\par
\R{直列}{ちょく|れつ}の\R{回路}{かい|ろ}で、\R{電源}{でん|げん}の\R{電圧}{でん|あつ}が、それぞれの\R{抵抗}{てい|こう}にかかる\R{電圧}{でん|あつ}の\R{和}{わ}になるのと\R{同}{おな}じです。\R{電源}{でん|げん}は\R{電気}{でん|き}を\R{高}{たか}い\R{所}{ところ}へもち\R{上}{あ}げるポンプ、\R{抵抗}{てい|こう}は\R{電気}{でん|き}が\R{下}{くだ}る\R{坂}{さか}のようなものです。\R{輪}{わ}をひとまわりするともとの\R{高}{たか}さにもどるので、もち\R{上}{あ}げた\R{高}{たか}さ（\R{電源}{でん|げん}の\R{電圧}{でん|あつ}）と、\R{下}{くだ}った\R{高}{たか}さの\R{合計}{ごう|けい}（\R{抵抗}{てい|こう}の\R{電圧}{でん|あつ}の\R{和}{わ}）は\R{等}{ひと}しくなります。

\R{左}{ひだり}の\R{輪}{わ}（\R{左}{ひだり}の\R{電源}{でん|げん}・\R{左}{ひだり}の\R{抵抗}{てい|こう}・まん\R{中}{なか}の\R{抵抗}{てい|こう}）と、\R{右}{みぎ}の\R{輪}{わ}（\R{右}{みぎ}の\R{電源}{でん|げん}・\R{右}{みぎ}の\R{抵抗}{てい|こう}・まん\R{中}{なか}の\R{抵抗}{てい|こう}）で、きまり3を\R{式}{しき}にします。
\begin{align*}
\text{\small\R{左}{ひだり}の\R{輪}{わ}}\quad 9&=2I_1+2(I_1+I_2) & &\to & 4I_1+2I_2&=9 && \cdots\text{①}\\
\text{\small\R{右}{みぎ}の\R{輪}{わ}}\quad 6&=2I_2+2(I_1+I_2) & &\to & 2I_1+4I_2&=6 && \cdots\text{②}
\end{align*}
\R{加減法}{か|げん|ほう}で\R{解}{と}きます。①を2\R{倍}{ばい}すると $8I_1+4I_2=18$。ここから②をひくと、$6I_1=12$、$I_1=2$。②に\R{入}{い}れて $4+4I_2=6$、$4I_2=2$、$I_2=0.5$。

\R{答}{こた}え：$I_1=\qty{2}{A}$、$I_2=\qty{0.5}{A}$、まん\R{中}{なか}の\R{抵抗}{てい|こう}には $2+0.5=2.5$ で\qty{2.5}{A}\R{流}{なが}れます。\R{確}{たし}かめ：\R{左}{ひだり}の\R{輪}{わ}は $2\times2+2\times2.5=4+5=9$、\R{右}{みぎ}の\R{輪}{わ}は $2\times0.5+2\times2.5=1+5=6$ で、どちらも\R{電源}{でん|げん}の\R{電圧}{でん|あつ}と\R{合}{あ}います。

2つの\R{電流}{でん|りゅう}は、どちらもまん\R{中}{なか}の\R{抵抗}{てい|こう}を\R{通}{とお}るので、\R{片方}{かた|ほう}だけでは\R{決}{き}まりません。2つの\R{輪}{わ}の\R{式}{しき}を\R{同時}{どう|じ}にみたす\R{組}{くみ}として、いっしょに\R{決}{き}まるのです。

\begin{motto}[3]
\textbf{3つの\R{文字}{も|じ}の\R{連立方程式}{れん|りつ|ほう|てい|しき}}\par
\R{輪}{わ}が3つある\R{回路}{かい|ろ}（\R{網目}{あみ|め}のような\R{回路}{かい|ろ}）では、わからない\R{電流}{でん|りゅう}が3つになり、\R{式}{しき}も3つ\R{必要}{ひつ|よう}です。たとえば、
\[ \begin{cases} x+y+z=9 & \cdots\text{①}\\ x-y=1 & \cdots\text{②}\\ y-z=1 & \cdots\text{③} \end{cases} \]
②から $x=y+1$、③から $z=y-1$。これを①に\R{代入}{だい|にゅう}すると、$(y+1)+y+(y-1)=9$、$3y=9$、$y=3$。だから $x=4$、$z=2$ です。\R{文字}{も|じ}を1つずつ\R{消}{け}していけば、\R{文字}{も|じ}がいくつあっても、\R{同}{おな}じ\R{考}{かんが}え\R{方}{かた}で\R{解}{と}けます。\R{式}{しき}の\R{数}{かず}が、わからない\R{数}{かず}の\R{数}{かず}と\R{同}{おな}じだけ\R{必要}{ひつ|よう}です。
\end{motto}

\begin{taisetsu}
\R{連立方程式}{れん|りつ|ほう|てい|しき}の\R{解}{かい}＝2つの\R{式}{しき}を\R{同時}{どう|じ}にみたす $(x,\ y)$ の\R{組}{くみ}＝2\R{本}{ほん}の\R{直線}{ちょく|せん}の\R{交点}{こう|てん}。\par\vspace{2pt}
\R{代入法}{だい|にゅう|ほう}：「$y=\cdots$」をもう1つの\R{式}{しき}にかっこごと\R{入}{い}れる。\R{加減法}{か|げん|ほう}：\R{係数}{けい|すう}をそろえて、\R{式}{しき}どうしをたす・ひく。\par\vspace{2pt}
\R{平行}{へい|こう}（\R{傾}{かたむ}きが\R{同}{おな}じ・\R{切片}{せっ|ぺん}がちがう）なら\R{解}{かい}なし。\R{重}{かさ}なるなら\R{解}{かい}は\R{無数}{む|すう}。\par\vspace{2pt}
\R{解}{かい}は、\R{必}{かなら}ず\R{両方}{りょう|ほう}の\R{式}{しき}に\R{代入}{だい|にゅう}して\R{確}{たし}かめる。
\end{taisetsu}

\begin{kotoba}
\small
\bkword{\R{連立方程式}{れん|りつ|ほう|てい|しき}}{2つ（いくつか）の\R{式}{しき}を\R{組}{くみ}にして、\R{同時}{どう|じ}に\R{成}{な}り\R{立}{た}たせる\R{文字}{も|じ}の\R{値}{あたい}をさがすもの。}
\bkword{\R{交点}{こう|てん}}{2\R{本}{ほん}の\R{線}{せん}が\R{交}{まじ}わる\R{点}{てん}。\R{両方}{りょう|ほう}の\R{式}{しき}をみたす。}
\bkword{\R{代入法}{だい|にゅう|ほう}}{1つの\R{式}{しき}を「$y=\cdots$」にして、もう1つの\R{式}{しき}に\R{代入}{だい|にゅう}し、\R{文字}{も|じ}を1つ\R{消}{け}す\R{解}{と}き\R{方}{かた}。}
\bkword{\R{加減法}{か|げん|ほう}}{\R{式}{しき}どうしをたしたりひいたりして、\R{文字}{も|じ}を1つ\R{消}{け}す\R{解}{と}き\R{方}{かた}。}
\bkword{\R{係数}{けい|すう}}{\R{文字}{も|じ}にかかっている\R{数}{かず}。$3y$ の3、$-y$ の $-1$。}
\end{kotoba}

%% ---------------- 例題 ----------------
\bksection{reidai}{\R{例題}{れい|だい}でたしかめよう}

\begin{reidai}[グラフで\R{交点}{こう|てん}を\R{読}{よ}み、\R{確}{たし}かめる]
\noindent\begin{minipage}[t]{0.54\linewidth}
$y=x+1$ と $y=-2x+7$ のグラフをかき、\R{交点}{こう|てん}を\R{読}{よ}みとりましょう。\R{読}{よ}んだ\R{点}{てん}が\R{両方}{りょう|ほう}の\R{式}{しき}をみたすか、\R{確}{たし}かめましょう。\par
\kangaekata \R{切片}{せっ|ぺん}と\R{傾}{かたむ}きからグラフをかきます（3-5）。$y=x+1$ は $(0,\ 1)$ から\R{右}{みぎ}へ1・\R{上}{うえ}へ1。$y=-2x+7$ は $(0,\ 7)$ から\R{右}{みぎ}へ1・\R{下}{した}へ2。\par
\shiki グラフから、\R{交点}{こう|てん}は $(2,\ 3)$\par
\hspace*{2.6\zw}$y=x+1$：$2+1=3$\par
\hspace*{2.6\zw}$y=-2x+7$：$-2\times2+7=3$\par
\kotae $(2,\ 3)$。\R{両方}{りょう|ほう}の\R{式}{しき}をみたす。
\end{minipage}\hfill
\begin{minipage}[t]{0.42\linewidth}\vspace{0pt}\centering
\begin{tikzpicture}
\begin{axis}[bkplot, width=52mm, height=54mm, xmin=-0.6, xmax=4.6, ymin=-0.6, ymax=7.8, xtick={1,...,4}, ytick={1,...,7}, xlabel={$x$}, ylabel={$y$}]
  \addplot[bkcurve, domain=-0.5:4.4] {x+1};
  \addplot[bkcurve2, domain=0:3.6] {-2*x+7};
  \draw[bkGray, line width=0.5pt, densely dashed] (axis cs:2,0) -- (axis cs:2,3) -- (axis cs:0,3);
  \addplot[only marks, mark=*, mark size=2pt, bkRed] coordinates {(2,3)};
  \node[anchor=west, font=\sffamily\scriptsize, text=bkBlue!80!black, inner sep=1pt] at (axis cs:3.2,5.6) {$y=x+1$};
  \node[anchor=west, font=\sffamily\scriptsize, text=bkOrange!80!black, inner sep=1pt] at (axis cs:1.1,6.3) {$y=-2x+7$};
\end{axis}
\end{tikzpicture}
\end{minipage}
\end{reidai}

\begin{jibun}
$y=x+2$ と $y=-x+6$ のグラフをノートにかき、\R{交点}{こう|てん}を\R{読}{よ}みとりましょう。\R{読}{よ}んだ\R{点}{てん}が\R{両方}{りょう|ほう}の\R{式}{しき}をみたすか、\R{確}{たし}かめましょう。
\bkjans{\R{交点}{こう|てん}は $(2,\ 4)$。$2+2=4$、$-2+6=4$ で\R{両方}{りょう|ほう}をみたす。}
\end{jibun}

\begin{reidai}[\R{代入法}{だい|にゅう|ほう}で\R{解}{と}く]
\R{次}{つぎ}の\R{連立方程式}{れん|りつ|ほう|てい|しき}を、\R{代入法}{だい|にゅう|ほう}で\R{解}{と}きましょう。\par
（1）\nobreak $\begin{cases} y=3x-4\\ y=-x+8 \end{cases}$\qquad（2）\nobreak $\begin{cases} 5x-2y=8\\ y=x-1 \end{cases}$\par
\kangaekata （1）は\R{両方}{りょう|ほう}とも「$y=$」なので、\R{右辺}{う|へん}どうしを＝で\R{結}{むす}びます。（2）は $2y$ の $y$ に $x-1$ をかっこごと\R{入}{い}れます。$-2\times(-1)=+2$ に\R{注意}{ちゅう|い}します。\par
\shiki（1）\nobreak $3x-4=-x+8$、$3x+x=8+4$、$4x=12$、$x=3$。$y=3\times3-4=5$\par
\hspace*{2.6\zw}\phantom{（1）}\R{確}{たし}かめ：$-3+8=5$\par
\hspace*{2.6\zw}（2）\nobreak $5x-2(x-1)=8$、$5x-2x+2=8$、$3x=6$、$x=2$。$y=2-1=1$\par
\hspace*{2.6\zw}\phantom{（2）}\R{確}{たし}かめ：$5\times2-2\times1=10-2=8$\par
\kotae（1）\nobreak $x=3$、$y=5$\quad（2）\nobreak $x=2$、$y=1$
\end{reidai}

\begin{reidai}[\R{加減法}{か|げん|ほう}：そのままひく・たす]
\R{次}{つぎ}の\R{連立方程式}{れん|りつ|ほう|てい|しき}を、\R{加減法}{か|げん|ほう}で\R{解}{と}きましょう。\par
（1）\nobreak $\begin{cases} 2x+3y=13\\ 2x+y=7 \end{cases}$\qquad（2）\nobreak $\begin{cases} x+y=7\\ x-y=1 \end{cases}$\par
\kangaekata （1）は $x$ の\R{係数}{けい|すう}が\R{同}{おな}じ2なので、ひくと $x$ が\R{消}{き}えます。（2）は $y$ の\R{係数}{けい|すう}が $+1$ と $-1$ なので、たすと $y$ が\R{消}{き}えます。\par
\shiki（1）\nobreak \R{上}{うえ}－\R{下}{した}：$3y-y=2y$、$13-7=6$。$2y=6$、$y=3$。$2x+3=7$、$2x=4$、$x=2$\par
\hspace*{2.6\zw}（2）\nobreak \R{上}{うえ}＋\R{下}{した}：$2x=8$、$x=4$。$4+y=7$、$y=3$\par
\hspace*{2.6\zw}\R{確}{たし}かめ：（1）$2\times2+3\times3=13$。（2）$4-3=1$\par
\kotae（1）\nobreak $x=2$、$y=3$\quad（2）\nobreak $x=4$、$y=3$
\end{reidai}

\begingroup\tcbset{equal height group=bkthrgjA}%
\noindent\begin{minipage}[t]{0.49\linewidth}
\begin{jibun}\raggedright
\R{次}{つぎ}の\R{連立方程式}{れん|りつ|ほう|てい|しき}を、\R{代入法}{だい|にゅう|ほう}で\R{解}{と}きましょう。（1）\nobreak $\begin{cases} y=2x+1\\ y=-x+7 \end{cases}$\quad（2）\nobreak $\begin{cases} 4x-3y=9\\ y=x-2 \end{cases}$
\bkjans{（1）\nobreak $2x+1=-x+7$、$3x=6$、$x=2$、$y=5$。（2）\nobreak $4x-3(x-2)=9$、$4x-3x+6=9$、$x=3$、$y=1$。}
\end{jibun}
\end{minipage}\hfill
\begin{minipage}[t]{0.49\linewidth}
\begin{jibun}\raggedright
\R{次}{つぎ}の\R{連立方程式}{れん|りつ|ほう|てい|しき}を、\R{加減法}{か|げん|ほう}で\R{解}{と}きましょう。（1）\nobreak $\begin{cases} 3x+2y=12\\ 3x+y=9 \end{cases}$\quad（2）\nobreak $\begin{cases} x+y=9\\ x-y=3 \end{cases}$
\bkjans{（1）\nobreak ひいて $y=3$。$3x+3=9$ より $x=2$。（2）\nobreak たして $2x=12$、$x=6$。$6+y=9$ より $y=3$。}
\end{jibun}
\end{minipage}\par
\endgroup

\begin{reidai}[\R{加減法}{か|げん|ほう}：かけてそろえてから]
\R{次}{つぎ}の\R{連立方程式}{れん|りつ|ほう|てい|しき}を、\R{加減法}{か|げん|ほう}で\R{解}{と}きましょう。\par
（1）\nobreak $\begin{cases} 2x+y=7 & \cdots\text{①}\\ x-3y=-7 & \cdots\text{②} \end{cases}$\qquad（2）\nobreak $\begin{cases} 2x+3y=12 & \cdots\text{①}\\ 3x+2y=13 & \cdots\text{②} \end{cases}$\par
\kangaekata （1）は①を3\R{倍}{ばい}すると、$y$ の\R{係数}{けい|すう}が $+3$ と $-3$ になるので、たします。（2）はどちらか1つだけではそろわないので、①を3\R{倍}{ばい}、②を2\R{倍}{ばい}して、$x$ の\R{係数}{けい|すう}を6にそろえ、ひきます。\R{右辺}{う|へん}にもかけるのを\R{忘}{わす}れないこと。\par
\shiki（1）\nobreak ①×3：$6x+3y=21$。これと②をたすと $7x=14$、$x=2$\par
\hspace*{2.6\zw}\phantom{（1）}①に\R{入}{い}れて $4+y=7$、$y=3$\par
\hspace*{2.6\zw}（2）\nobreak ①×3：$6x+9y=36$。②×2：$6x+4y=26$\par
\hspace*{2.6\zw}\phantom{（2）}ひくと $5y=10$、$y=2$\par
\hspace*{2.6\zw}\phantom{（2）}①に\R{入}{い}れて $2x+6=12$、$2x=6$、$x=3$\par
\hspace*{2.6\zw}\R{確}{たし}かめ：（1）②は $2-9=-7$。（2）②は $9+4=13$\par
\kotae（1）\nobreak $x=2$、$y=3$\quad（2）\nobreak $x=3$、$y=2$
\end{reidai}

\begin{jibun}
\R{次}{つぎ}の\R{連立方程式}{れん|りつ|ほう|てい|しき}を、\R{加減法}{か|げん|ほう}で\R{解}{と}きましょう。\par
（1）\nobreak $\begin{cases} x+3y=10\\ 2x-y=6 \end{cases}$\par
（2）\nobreak $\begin{cases} 3x+2y=7\\ 2x+3y=8 \end{cases}$
\bkjans{（1）\nobreak \R{下}{した}の\R{式}{しき}×3 は $6x-3y=18$。\R{上}{うえ}とたして $7x=28$、$x=4$。$4+3y=10$ より $y=2$。（2）\nobreak \R{上}{うえ}×2 は $6x+4y=14$、\R{下}{した}×3 は $6x+9y=24$。ひいて $5y=10$、$y=2$。$3x+4=7$ より $x=1$。}
\end{jibun}

\begin{reidai}[\R{買}{か}い\R{物}{もの}の\R{問題}{もん|だい}]
えんぴつ3\R{本}{ぼん}とノート2\R{冊}{さつ}で520\R{円}{えん}、えんぴつ1\R{本}{ぽん}とノート3\R{冊}{さつ}で570\R{円}{えん}でした。えんぴつ1\R{本}{ぽん}とノート1\R{冊}{さつ}の\R{値段}{ね|だん}を\R{求}{もと}めましょう。\par
\kangaekata わからない\R{数}{かず}が2つなので、2つの\R{文字}{も|じ}にします。えんぴつ1\R{本}{ぽん}を $x$ \R{円}{えん}、ノート1\R{冊}{さつ}を $y$ \R{円}{えん}とします。\R{代金}{だい|きん}＝\R{値段}{ね|だん}×\R{数}{かず}の\R{合計}{ごう|けい}で、2つの\R{式}{しき}を\R{作}{つく}ります。\par
\shiki $3x+2y=520$\ ①\qquad $x+3y=570$\ ②\par
\hspace*{2.6\zw}②×3：$3x+9y=1710$。これから①をひくと $7y=1190$、$y=170$\par
\hspace*{2.6\zw}②に\R{入}{い}れて $x+510=570$、$x=60$\par
\hspace*{2.6\zw}\R{確}{たし}かめ：①は $180+340=520$\par
\kotae えんぴつ60\R{円}{えん}、ノート170\R{円}{えん}
\end{reidai}

\begin{reidai}[\R{出会}{で|あ}いの\R{問題}{もん|だい}]
\R{家}{いえ}と\R{駅}{えき}は\qty{1200}{m}はなれています。Aさんは\R{家}{いえ}から\R{分速}{ふん|そく}\qty{50}{m}で、Bさんは\R{同}{おな}じときに\R{駅}{えき}から\R{分速}{ふん|そく}\qty{70}{m}で、おたがいに\R{向}{む}かって\R{歩}{ある}きます。\R{出会}{で|あ}うのは\R{何分後}{なん|ぷん|ご}、\R{家}{いえ}から\R{何}{なん}mの\R{所}{ところ}ですか。\par
\kangaekata $x$ \R{分後}{ふん|ご}の\R{家}{いえ}からの\R{位置}{い|ち}を $y$〔m〕として、2\R{人}{にん}の\R{式}{しき}を\R{作}{つく}ります。Bさんは1200から $70x$ ずつ\R{小}{ちい}さくなります。\par
\shiki A：$y=50x$\quad B：$y=-70x+1200$\par
\hspace*{2.6\zw}$50x=-70x+1200$、$120x=1200$、$x=10$。$y=50\times10=500$\par
\hspace*{2.6\zw}\R{確}{たし}かめ：$-700+1200=500$\par
\kotae 10\R{分後}{ぷん|ご}、\R{家}{いえ}から\qty{500}{m}の\R{所}{ところ}
\end{reidai}

\begingroup\tcbset{equal height group=bkthrgjB}%
\noindent\begin{minipage}[t]{0.49\linewidth}
\begin{jibun}\raggedright
ジュース2\R{本}{ほん}とパン1\R{個}{こ}で380\R{円}{えん}、ジュース1\R{本}{ぽん}とパン3\R{個}{こ}で540\R{円}{えん}でした。ジュース1\R{本}{ぽん}とパン1\R{個}{こ}の\R{値段}{ね|だん}を\R{求}{もと}めましょう。
\bkjans{ジュースを $x$ \R{円}{えん}、パンを $y$ \R{円}{えん}とすると、$2x+y=380$、$x+3y=540$。2つめ×2 は $2x+6y=1080$。1つめをひいて $5y=700$、$y=140$。$x=540-420=120$。ジュース120\R{円}{えん}、パン140\R{円}{えん}。}
\end{jibun}
\end{minipage}\hfill
\begin{minipage}[t]{0.49\linewidth}
\begin{jibun}\raggedright
\R{家}{いえ}と\R{駅}{えき}は\qty{1500}{m}はなれています。Cさんは\R{家}{いえ}から\R{分速}{ふん|そく}\qty{40}{m}で、Dさんは\R{同}{おな}じときに\R{駅}{えき}から\R{分速}{ふん|そく}\qty{60}{m}で、おたがいに\R{向}{む}かって\R{歩}{ある}きます。\R{出会}{で|あ}うのは\R{何分後}{なん|ぷん|ご}、\R{家}{いえ}から\R{何}{なん}mの\R{所}{ところ}ですか。
\bkjans{$y=40x$、$y=-60x+1500$。$40x=-60x+1500$、$100x=1500$、$x=15$。$y=40\times15=600$。15\R{分後}{ふん|ご}、\R{家}{いえ}から\qty{600}{m}の\R{所}{ところ}。}
\end{jibun}
\end{minipage}\par
\endgroup

\begin{reidai}[\R{解}{かい}はいくつ？]
\R{次}{つぎ}の\R{連立方程式}{れん|りつ|ほう|てい|しき}の\R{解}{かい}は、1\R{組}{くみ}・なし・\R{無数}{む|すう}のどれですか。1\R{組}{くみ}なら\R{求}{もと}めましょう。\par
（1）\nobreak $\begin{cases} y=2x+1\\ y=2x-3 \end{cases}$\quad（2）\nobreak $\begin{cases} y=-x+4\\ 2x+2y=8 \end{cases}$\quad（3）\nobreak $\begin{cases} y=2x+1\\ y=-x+4 \end{cases}$\par
\kangaekata 2つの\R{式}{しき}を「$y=ax+b$」の\R{形}{かたち}にして、\R{傾}{かたむ}きと\R{切片}{せっ|ぺん}をくらべます。\par
\shiki（1）\nobreak \R{傾}{かたむ}きはどちらも2、\R{切片}{せっ|ぺん}は1と $-3$ でちがう、\R{平行}{へい|こう}\par
\hspace*{2.6\zw}（2）\nobreak $2y=-2x+8$、$y=-x+4$。1つめと\R{同}{おな}じ\R{直線}{ちょく|せん}、\R{重}{かさ}なる\par
\hspace*{2.6\zw}（3）\nobreak \R{傾}{かたむ}きが2と $-1$ でちがう、1\R{点}{てん}で\R{交}{まじ}わる。$2x+1=-x+4$、$3x=3$、$x=1$、$y=3$\par
\kotae（1）\nobreak なし\quad（2）\nobreak \R{無数}{む|すう}\quad（3）\nobreak 1\R{組}{くみ}、$x=1$、$y=3$
\end{reidai}

\begin{reidai}[2\R{点}{てん}を\R{通}{とお}る\R{直線}{ちょく|せん}を、\R{連立方程式}{れん|りつ|ほう|てい|しき}で\R{求}{もと}める]
2\R{点}{てん} $(1,\ 5)$、$(3,\ 11)$ を\R{通}{とお}る\R{直線}{ちょく|せん} $y=ax+b$ の $a$ と $b$ を、\R{連立方程式}{れん|りつ|ほう|てい|しき}で\R{求}{もと}めましょう。\par
\kangaekata 3-5では\R{傾}{かたむ}きを\R{先}{さき}に\R{求}{もと}めました。ここでは、わからない $a$ と $b$ を2つの\R{文字}{も|じ}と\R{考}{かんが}えます。\R{直線}{ちょく|せん}が\R{点}{てん}を\R{通}{とお}るので、その $x$、$y$ を\R{代入}{だい|にゅう}すると\R{等式}{とう|しき}が\R{成}{な}り\R{立}{た}ちます。\par
\shiki $(1,\ 5)$ を\R{代入}{だい|にゅう}：$5=a\times1+b$、$a+b=5$\ ①\par
\hspace*{2.6\zw}$(3,\ 11)$ を\R{代入}{だい|にゅう}：$11=a\times3+b$、$3a+b=11$\ ②\par
\hspace*{2.6\zw}②－①：$2a=6$、$a=3$。①に\R{入}{い}れて $3+b=5$、$b=2$\par
\kotae $a=3$、$b=2$（$y=3x+2$。3-5の\R{例題}{れい|だい}と\R{同}{おな}じ\R{答}{こた}え）
\end{reidai}

\begingroup\tcbset{equal height group=bkthrgjC}%
\noindent\begin{minipage}[t]{0.49\linewidth}
\begin{jibun}\raggedright
\R{次}{つぎ}の\R{連立方程式}{れん|りつ|ほう|てい|しき}の\R{解}{かい}は、1\R{組}{くみ}・なし・\R{無数}{む|すう}のどれですか。1\R{組}{くみ}なら\R{求}{もと}めましょう。（1）\nobreak $\begin{cases} y=3x-1\\ y=3x+2 \end{cases}$\quad（2）\nobreak $\begin{cases} y=x+3\\ 3x-3y=-9 \end{cases}$\quad（3）\nobreak $\begin{cases} y=x+3\\ y=-2x+6 \end{cases}$
\bkjans{（1）\nobreak \R{傾}{かたむ}きが\R{同}{おな}じで\R{切片}{せっ|ぺん}がちがうので、なし。（2）\nobreak $-3y=-3x-9$、$y=x+3$ で\R{同}{おな}じ\R{直線}{ちょく|せん}なので、\R{無数}{む|すう}。（3）\nobreak $x+3=-2x+6$、$3x=3$、$x=1$、$y=4$。1\R{組}{くみ}。}
\end{jibun}
\end{minipage}\hfill
\begin{minipage}[t]{0.49\linewidth}
\begin{jibun}\raggedright
2\R{点}{てん} $(2,\ 7)$、$(4,\ 11)$ を\R{通}{とお}る\R{直線}{ちょく|せん} $y=ax+b$ の $a$ と $b$ を、\R{連立方程式}{れん|りつ|ほう|てい|しき}で\R{求}{もと}めましょう。
\bkjans{$2a+b=7$、$4a+b=11$。ひいて $2a=4$、$a=2$。$4+b=7$ より $b=3$。$y=2x+3$。}
\end{jibun}
\end{minipage}\par
\endgroup

\begin{reidai}[\R{回路}{かい|ろ}の\R{電流}{でん|りゅう}を\R{求}{もと}める]
\noindent\begin{minipage}[t]{0.34\linewidth}
\R{右}{みぎ}の\R{回路}{かい|ろ}で、\R{電流}{でん|りゅう} $I_1$、$I_2$ と、まん\R{中}{なか}の\R{抵抗}{てい|こう}を\R{流}{なが}れる\R{電流}{でん|りゅう}を\R{求}{もと}めましょう。
\end{minipage}\hfill
\begin{minipage}[t]{0.64\linewidth}\vspace{0pt}\centering
\begin{tikzpicture}[x=0.62mm, y=0.62mm]
  \bkthrgcircuit{\qty{6}{V}}{\qty{3}{\ohm}}{\qty{1}{\ohm}}{\qty{3}{\ohm}}{\qty{9}{V}}
\end{tikzpicture}
\end{minipage}\par
\kangaekata まん\R{中}{なか}の\R{電流}{でん|りゅう}は $I_1+I_2$（きまり1）。\R{抵抗}{てい|こう}の\R{電圧}{でん|あつ}は $RI$（きまり2）。\R{左}{ひだり}の\R{輪}{わ}と\R{右}{みぎ}の\R{輪}{わ}で「\R{電源}{でん|げん}の\R{電圧}{でん|あつ}＝\R{抵抗}{てい|こう}の\R{電圧}{でん|あつ}の\R{和}{わ}」（きまり3）。\par
\shiki \R{左}{ひだり}：$6=3I_1+1\times(I_1+I_2)$、$4I_1+I_2=6$\ ①\par
\hspace*{2.6\zw}\R{右}{みぎ}：$9=3I_2+1\times(I_1+I_2)$、$I_1+4I_2=9$\ ②\par
\hspace*{2.6\zw}①×4：$16I_1+4I_2=24$。これから②をひくと $15I_1=15$、$I_1=1$\par
\hspace*{2.6\zw}①に\R{入}{い}れて $4+I_2=6$、$I_2=2$。まん\R{中}{なか}は $1+2=3$\par
\hspace*{2.6\zw}\R{確}{たし}かめ：\R{左}{ひだり} $3\times1+1\times3=6$、\R{右}{みぎ} $3\times2+1\times3=9$\par
\kotae $I_1=\qty{1}{A}$、$I_2=\qty{2}{A}$、まん\R{中}{なか}は\qty{3}{A}
\end{reidai}

\begin{jibun}
\bkreidairef{9}の\R{回路}{かい|ろ}で、\R{左}{ひだり}の\R{電源}{でん|げん}を\qty{9}{V}、\R{右}{みぎ}の\R{電源}{でん|げん}を\qty{6}{V}にかえました（\R{抵抗}{てい|こう}は\R{左}{ひだり}\qty{3}{\ohm}、まん\R{中}{なか}\qty{1}{\ohm}、\R{右}{みぎ}\qty{3}{\ohm}のまま）。$I_1$、$I_2$ と、まん\R{中}{なか}の\R{電流}{でん|りゅう}を\R{求}{もと}めましょう。
\bkjans{\R{左}{ひだり}：$9=3I_1+(I_1+I_2)$ より $4I_1+I_2=9$。\R{右}{みぎ}：$6=3I_2+(I_1+I_2)$ より $I_1+4I_2=6$。1つめ×4 から2つめをひいて $15I_1=30$、$I_1=2$。$8+I_2=9$ より $I_2=1$。$I_1=\qty{2}{A}$、$I_2=\qty{1}{A}$、まん\R{中}{なか}は\qty{3}{A}。}
\end{jibun}

%% ---------------- まちがえやすいところ ----------------
\bksection{tsumazuki}{まちがえやすいところをチェック}\label{bk:3-6:tsuma}

\begin{tsumazuki}{\R{代入}{だい|にゅう}するとき、かっこをつけない}
\bkng{$3x+2y=12$ に $y=2x-1$ を\R{入}{い}れて、$3x+2\times2x-1=12$}
\bkok{$3x+2(2x-1)=12$ → $3x+4x-2=12$}
×では、2を $2x$ にしかかけていません。$-1$ にもかける\R{必要}{ひつ|よう}があります。$y$ の\R{箱}{はこ}に、$2x-1$ をまるごと\R{入}{い}れよう。
\end{tsumazuki}

\begin{tsumazuki}{ひくときの\R{符号}{ふ|ごう}}
\bkng{$2x+3y=13$ から $2x-y=1$ をひいて、$2y=12$}
\bkok{$3y-(-y)=3y+y=4y$。$4y=12$ → $y=3$}
マイナスの\R{項}{こう}をひくと、プラスになります（2-2）。\R{心配}{しん|ぱい}なら、\R{下}{した}の\R{式}{しき}の\R{符号}{ふ|ごう}を\R{全部}{ぜん|ぶ}\R{反対}{はん|たい}にしてから、たすとよい。
\end{tsumazuki}

\begin{tsumazuki}{\R{右辺}{う|へん}にかけ\R{忘}{わす}れる}
\bkng{$x+2y=8$ を2\R{倍}{ばい}して、$2x+4y=8$}
\bkok{$2x+4y=16$（\R{右辺}{う|へん}の8にも2をかける）}
\R{天}{てん}びんの\R{片方}{かた|ほう}の\R{皿}{さら}だけ2\R{倍}{ばい}にすると、つり\R{合}{あ}いがくずれます。\R{両辺}{りょう|へん}の\R{全部}{ぜん|ぶ}の\R{項}{こう}にかけます。
\end{tsumazuki}

\begin{tsumazuki}{\R{片方}{かた|ほう}だけ\R{求}{もと}めて\R{終}{お}わる}
\bkng{$y=3$ が\R{出}{で}たので、\R{答}{こた}えは $y=3$}
\bkok{$x$ も\R{求}{もと}めて「$x=2$、$y=3$」}
\R{連立方程式}{れん|りつ|ほう|てい|しき}の\R{解}{かい}は、$x$ と $y$ の\R{組}{くみ}です。\R{交点}{こう|てん}の\R{場所}{ば|しょ}は、2つの\R{数}{かず}がそろわないと\R{決}{き}まりません。
\end{tsumazuki}

\begin{tsumazuki}{$0=10$ になったら「\R{解}{かい}は0」？}
\bkng{$2x=2x+10$ を\R{解}{と}くと $0=10$。だから $x=0$}
\bkok{どんな $x$ でも\R{成}{な}り\R{立}{た}たないので、\R{解}{かい}なし}
$x=0$ を\R{入}{い}れると $0=10$ で、\R{成}{な}り\R{立}{た}ちません。\R{傾}{かたむ}きが\R{同}{おな}じ2\R{本}{ほん}の\R{平行}{へい|こう}な\R{線}{せん}なので、\R{交点}{こう|てん}がないのです。
\end{tsumazuki}

\begin{column}{「\R{方程式}{ほう|てい|しき}」ということばは、\R{連立方程式}{れん|りつ|ほう|てい|しき}から}
\R{今}{いま}から2000\R{年}{ねん}ほど\R{前}{まえ}の\R{中国}{ちゅう|ごく}の\R{数学}{すう|がく}の\R{本}{ほん}『\R{九章算術}{きゅう|しょう|さん|じゅつ}』には、「\R{方程}{ほう|てい}」という\R{章}{しょう}があります。そこでは、\R{何種類}{なん|しゅ|るい}かの\R{穀物}{こく|もつ}の\R{量}{りょう}と\R{合計}{ごう|けい}を、\R{算木}{さん|ぎ}という\R{細}{ほそ}い\R{棒}{ぼう}で\R{表}{ひょう}のようにならべ、\R{列}{れつ}どうしをひいたり\R{何倍}{なん|ばい}かしたりして、\R{連立方程式}{れん|りつ|ほう|てい|しき}を\R{解}{と}いていました。⑥の「\R{係数}{けい|すう}を\R{表}{ひょう}にならべる」\R{方法}{ほう|ほう}と、ほとんど\R{同}{おな}じです。「\R{方程式}{ほう|てい|しき}」ということばは、この「\R{方程}{ほう|てい}」から\R{来}{き}ています。
\end{column}

%% ---------------- 物理・工学 ----------------
\begin{butsuri}
\par\addvspace{3pt}\textbf{2\R{人}{にん}で\R{荷物}{に|もつ}をかつぐ：\R{力}{ちから}のつり\R{合}{あ}い}\par\noindent
\begin{minipage}[c]{0.56\linewidth}
\R{長}{なが}さ\qty{3}{m}の\R{棒}{ぼう}の、Aさんから\qty{1}{m}の\R{所}{ところ}に、\R{重}{おも}さ\qty{300}{N}（\R{質量}{しつ|りょう}\R{約}{やく}\qty{30}{kg}の\R{物}{もの}にはたらく\R{重力}{じゅう|りょく}。2-15）の\R{荷物}{に|もつ}をつるし、AさんとBさんが\R{両}{りょう}はしでかつぎます。2\R{人}{にん}が\R{棒}{ぼう}を\R{上}{うえ}へおす\R{力}{ちから}を $F_A$、$F_B$（N）とします。\R{棒}{ぼう}は\R{動}{うご}かないので、
\end{minipage}\hfill
\begin{minipage}[c]{0.4\linewidth}\centering
\begin{tikzpicture}[x=0.85mm, y=0.85mm]
  \fill[bkStem!70, rounded corners=0.6mm] (-1,-1) rectangle (49,1);
  \draw[bkInk!70, line width=0.5pt] (16,-1) -- (16,-5);
  \filldraw[fill=bkGray!45, draw=bkGray!90!black, line width=0.6pt, rounded corners=0.8mm] (12,-12) rectangle (20,-5);
  \node[font=\sffamily\scriptsize] at (16,-8.5) {\qty{300}{N}};
  \draw[bkRed, line width=1pt, -{Stealth[length=4pt]}] (0,-9) -- (0,-1.8);
  \draw[bkRed, line width=1pt, -{Stealth[length=4pt]}] (48,-9) -- (48,-1.8);
  \node[font=\sffamily\scriptsize, text=bkRed!80!black, anchor=north] at (0,-9.5) {$F_A$};
  \node[font=\sffamily\scriptsize, text=bkRed!80!black, anchor=north] at (48,-9.5) {$F_B$};
  \draw[{Stealth[length=3pt]}-{Stealth[length=3pt]}, bkBlue, line width=0.5pt] (0,3.5) -- (16,3.5);
  \draw[{Stealth[length=3pt]}-{Stealth[length=3pt]}, bkBlue, line width=0.5pt] (16,3.5) -- (48,3.5);
  \node[font=\sffamily\scriptsize, text=bkBlue!80!black, anchor=south] at (8,4) {\qty{1}{m}};
  \node[font=\sffamily\scriptsize, text=bkBlue!80!black, anchor=south] at (32,4) {\qty{2}{m}};
  \node[font=\sffamily\scriptsize, anchor=south] at (0,8.5) {A};
  \node[font=\sffamily\scriptsize, anchor=south] at (48,8.5) {B};
\end{tikzpicture}
\end{minipage}
\par\vspace{2pt}
① \R{上}{うえ}と\R{下}{した}の\R{力}{ちから}がつり\R{合}{あ}う：$F_A+F_B=300$\par
② \R{荷物}{に|もつ}の\R{所}{ところ}を\R{支点}{し|てん}と\R{見}{み}たてこのつり\R{合}{あ}い（\R{小}{しょう}6の\R{理科}{り|か}・2-15）：$F_A\times1=F_B\times2$\par
②から $F_A=2F_B$。①に\R{代入}{だい|にゅう}して $3F_B=300$、$F_B=100$、$F_A=200$。\R{荷物}{に|もつ}に\R{近}{ちか}いAさんが、Bさんの2\R{倍}{ばい}の\qty{200}{N}を\R{支}{ささ}えます。\R{橋}{はし}をささえる\R{両}{りょう}はしの\R{柱}{はしら}にかかる\R{力}{ちから}も、\R{建築}{けん|ちく}ではこの\R{連立方程式}{れん|りつ|ほう|てい|しき}で\R{計算}{けい|さん}します。

\par\addvspace{3pt}\textbf{スマートフォンの\R{位置}{い|ち}を\R{決}{き}める}\par\noindent
スマートフォンは、\R{空}{そら}の4つ\R{以上}{い|じょう}の\R{人工衛星}{じん|こう|えい|せい}から\R{届}{とど}く\R{電波}{でん|ぱ}をもとに、\R{自分}{じ|ぶん}の\R{位置}{い|ち}を\R{計算}{けい|さん}しています（GPS）。わからない\R{数}{かず}は、\R{東西}{とう|ざい}・\R{南北}{なん|ぼく}・\R{高}{たか}さの3つと、スマートフォンの\R{時計}{と|けい}のずれの、\R{合}{あ}わせて4つ。だから、4つの\R{衛星}{えい|せい}から4つの\R{式}{しき}を\R{作}{つく}り、\R{連立方程式}{れん|りつ|ほう|てい|しき}として\R{同時}{どう|じ}に\R{解}{と}いています。わからない\R{数}{かず}と\R{同}{おな}じ\R{数}{かず}だけ、\R{式}{しき}が\R{必要}{ひつ|よう}なのです。
\end{butsuri}

\begin{column}{\R{工場}{こう|じょう}で\R{液}{えき}を\R{混}{ま}ぜる}
\R{化学}{か|がく}\R{工場}{こう|じょう}では、こさのちがう\R{液}{えき}を\R{混}{ま}ぜて、ちょうどよいこさの\R{液}{えき}を\R{作}{つく}ります。10\%の\R{食塩水}{しょく|えん|すい}と4\%の\R{食塩水}{しょく|えん|すい}を\R{混}{ま}ぜて、6\%の\R{食塩水}{しょく|えん|すい}を\qty{300}{g}\R{作}{つく}るとします。10\%のほうを $x$〔g〕、4\%のほうを $y$〔g〕とすると、\R{全体}{ぜん|たい}の\R{重}{おも}さで $x+y=300$、とけている\R{食塩}{しょく|えん}の\R{重}{おも}さで $0.1x+0.04y=300\times0.06=18$。2つめを100\R{倍}{ばい}すると $10x+4y=1800$。1つめを4\R{倍}{ばい}した $4x+4y=1200$ をひいて、$6x=600$、$x=100$、$y=200$ です。
\end{column}



%% ---------------- 章末問題 ----------------
\bkendproblems
\bkproblemgroup{\R{基本}{き|ほん}}{2}

\bkproblem $y=x+3$ と $y=-x+5$ のグラフをノートにかき、\R{交点}{こう|てん}を\R{読}{よ}みとりましょう。\R{両方}{りょう|ほう}の\R{式}{しき}をみたすかも\R{確}{たし}かめましょう。
\bkpans{$(1,\ 4)$。$1+3=4$、$-1+5=4$。}

\bkproblem \R{代入法}{だい|にゅう|ほう}で\R{解}{と}きましょう。（1）\nobreak $\begin{cases} y=2x-3\\ y=-x+6 \end{cases}$\quad（2）\nobreak $\begin{cases} 2x-3y=1\\ y=x-1 \end{cases}$
\bkpans{（1）\nobreak $2x-3=-x+6$、$3x=9$、$x=3$、$y=3$。（2）\nobreak $2x-3(x-1)=1$、$-x+3=1$、$x=2$、$y=1$。}

\bkproblem \R{加減法}{か|げん|ほう}で\R{解}{と}きましょう。（1）\nobreak $\begin{cases} 4x+3y=18\\ 4x+y=14 \end{cases}$\quad（2）\nobreak $\begin{cases} x+y=10\\ x-y=4 \end{cases}$
\bkpans{（1）\nobreak ひいて $2y=4$、$y=2$。$4x+2=14$ より $x=3$。（2）\nobreak たして $2x=14$、$x=7$。$y=3$。}

\bkproblem \R{加減法}{か|げん|ほう}で\R{解}{と}きましょう。（1）\nobreak $\begin{cases} 3x+y=10\\ x-2y=1 \end{cases}$\quad（2）\nobreak $\begin{cases} 2x+3y=7\\ 3x+2y=8 \end{cases}$
\bkpans{（1）\nobreak \R{上}{うえ}×2 は $6x+2y=20$。\R{下}{した}とたして $7x=21$、$x=3$。$9+y=10$ より $y=1$。（2）\nobreak \R{上}{うえ}×3 は $6x+9y=21$、\R{下}{した}×2 は $6x+4y=16$。ひいて $5y=5$、$y=1$。$2x+3=7$ より $x=2$。}

\bkproblem \R{解}{かい}は1\R{組}{くみ}・なし・\R{無数}{む|すう}のどれですか。\par
\hspace*{1\zw}（1）\nobreak $\begin{cases} y=-2x+5\\ y=-2x+1 \end{cases}$\qquad（2）\nobreak $\begin{cases} y=x-2\\ 2x-2y=4 \end{cases}$
\bkpans{（1）\nobreak \R{傾}{かたむ}きが\R{同}{おな}じで\R{切片}{せっ|ぺん}がちがうので、なし。（2）\nobreak $-2y=-2x+4$、$y=x-2$ で\R{同}{おな}じ\R{直線}{ちょく|せん}なので、\R{無数}{む|すう}。}

\bkproblem \R{博物館}{はく|ぶつ|かん}の\R{入館料}{にゅう|かん|りょう}は、\R{大人}{おとな}2\R{人}{にん}と\R{子}{こ}ども3\R{人}{にん}で2600\R{円}{えん}、\R{大人}{おとな}1\R{人}{にん}と\R{子}{こ}ども2\R{人}{にん}で1500\R{円}{えん}でした。\R{大人}{おとな}1\R{人}{にん}と\R{子}{こ}ども1\R{人}{にん}の\R{入館料}{にゅう|かん|りょう}を\R{求}{もと}めましょう。
\bkpans{\R{大人}{おとな}を $x$ \R{円}{えん}、\R{子}{こ}どもを $y$ \R{円}{えん}とすると、$2x+3y=2600$、$x+2y=1500$。2つめ×2 は $2x+4y=3000$。1つめをひいて $y=400$。$x=1500-800=700$。\R{大人}{おとな}700\R{円}{えん}、\R{子}{こ}ども400\R{円}{えん}。}

\bkproblem \qty{2000}{m}はなれた2つの\R{地点}{ち|てん}から、Eさんは\R{分速}{ふん|そく}\qty{60}{m}、Fさんは\R{分速}{ふん|そく}\qty{100}{m}で、\R{同}{おな}じときにおたがいに\R{向}{む}かって\R{歩}{ある}きます。\R{出会}{で|あ}うのは\R{何分後}{なん|ぷん|ご}、Eさんの\R{出発}{しゅっ|ぱつ}した\R{地点}{ち|てん}から\R{何}{なん}mの\R{所}{ところ}ですか。
\bkpans{$y=60x$、$y=-100x+2000$。$160x=2000$、$x=12.5$。$y=60\times12.5=750$。12.5\R{分後}{ふん|ご}（12\R{分}{ふん}30\R{秒後}{びょう|ご}）、\qty{750}{m}の\R{所}{ところ}。}

\bkproblemgroup{チャレンジ}{3}

\bkproblem \noindent\begin{minipage}[t]{0.3\linewidth}
\R{右}{みぎ}の\R{回路}{かい|ろ}で、\R{電流}{でん|りゅう} $I_1$、$I_2$ と、まん\R{中}{なか}の\R{抵抗}{てい|こう}を\R{流}{なが}れる\R{電流}{でん|りゅう}を\R{求}{もと}めましょう。
\end{minipage}\hfill
\begin{minipage}[t]{0.62\linewidth}\vspace{0pt}\centering
\begin{tikzpicture}[x=0.57mm, y=0.57mm]
  \bkthrgcircuit{\qty{9}{V}}{\qty{3}{\ohm}}{\qty{2}{\ohm}}{\qty{3}{\ohm}}{\qty{12}{V}}
\end{tikzpicture}
\end{minipage}
\bkpans{\R{左}{ひだり}：$9=3I_1+2(I_1+I_2)$ より $5I_1+2I_2=9$。\R{右}{みぎ}：$12=3I_2+2(I_1+I_2)$ より $2I_1+5I_2=12$。1つめ×5 は $25I_1+10I_2=45$、2つめ×2 は $4I_1+10I_2=24$。ひいて $21I_1=21$、$I_1=1$。$5+2I_2=9$ より $I_2=2$。$I_1=\qty{1}{A}$、$I_2=\qty{2}{A}$、まん\R{中}{なか}は\qty{3}{A}。}

\bkproblem 2\R{点}{てん} $(-1,\ 1)$、$(2,\ 7)$ を\R{通}{とお}る\R{直線}{ちょく|せん} $y=ax+b$ の $a$ と $b$ を、\R{連立方程式}{れん|りつ|ほう|てい|しき}で\R{求}{もと}めましょう。
\bkpans{$-a+b=1$、$2a+b=7$。2つめから1つめをひいて $3a=6$、$a=2$。$-2+b=1$ より $b=3$。$y=2x+3$。}

\begin{furikaeri}
\begin{checklist}
\item \R{交点}{こう|てん}が、2つの\R{式}{しき}を\R{同時}{どう|じ}にみたす\R{組}{くみ}だとわかる。\bkfurref{bk:3-6:kouten}
\item \R{代入法}{だい|にゅう|ほう}と\R{加減法}{か|げん|ほう}で\R{解}{と}き、\R{両方}{りょう|ほう}の\R{式}{しき}で\R{確}{たし}かめられる。\bkfurref{bk:3-6:dainyuu}
\item \R{解}{かい}がない・\R{無数}{む|すう}のわけを\R{説明}{せつ|めい}し、\R{出会}{で|あ}い・\R{買}{か}い\R{物}{もの}・\R{回路}{かい|ろ}の\R{問題}{もん|だい}を\R{解}{と}ける。\bkfurref{bk:3-6:heikou}
\end{checklist}
\end{furikaeri}

\bkfukidashi{mina}{2つのわからない\R{数}{かず}は、2つの\R{式}{しき}でいっしょに\R{決}{き}まるんだね。\R{次}{つぎ}の\R{部}{ぶ}では、\R{曲}{ま}がったグラフを\R{調}{しら}べるよ。}
